\bookStart{Speeches of Grimner}[Grímnis mǫ́l]
\setBookCode{Grimnismal}

\begin{flushright}%
\textbf{Dating} \parencite{Sapp2022}: C10th (0.976)

\textbf{Meter:} \Ljodahattr\ (1--2/2, 3--26, 27/4--27/7, 28/1--28/2, 28/6--28/7, 29--33/2, 35--45/2, 46/1--46/2, 47--48/2, 49/3--52, 54--57), \Fornyrdislag\ (2/3--2/4, 28/3--28/5, 33/3, 45/3--45/5, 48/3--48/4, 49/1--49/2, 53), \Galdralag\ (27/1--27/3, 34, 46/3--46/5)%
\end{flushright}

\section{Introduction}

The \textbf{Speeches of Grimner} (abbrev.~\textlink{Grimnismal}) is a mythological gnomic poem consisting of Weden’s wisdom told to the young prince Ayner as the god sits bound between two fires.

\subsection{Textual preservation and transmission}

\Grimnismal\ is preserved whole in both \Regius\ and \AM.  The shared prose passages between them prove beyond any doubt that the version of \Grimnismal\ in these two manuscripts ultimately goes back to a single written archetype, which may be called *\Regius\AM.

22 stanzas (11--15, 18--20, 23--24, 29, 34--37, 41--42, 45, 47--49) of \Grimnismal\ are preserved in citation in \Gylfaginning, in the case of 23, 24, and 37 explicitly attributed to \emph{Grímnis mǫ́l}.  Numerous other stanzas (6--8, TODO) are paraphrased in \Gylfaginning, more or less directly.  These stanzas may be said to belong to a \GylfMS\ redaction of the poem.

\Gylfaginning\ predates both \Regius\ and \AM, and \GylfMS\ must thus be distinct from both of them.  Common innovations shared between \Regius\ and \AM\ (e.g.~st.~18/4: om.~\emph{við}, 45/4: \emph{Bil·rǫst}) show that the two share a common written archetype, distinct from that of \GylfMS.  It is unclear whether \GylfMS\ reflects a more distantly related written version or an alternative oral redaction of the poem.

\subsection{Structure}

\textlink{Grimnismal} essentially consists of several nested layers.  The outermost layer is the prose passages which bracket the actual poem (P1--P2, P3).  It is hard to say for how long these have accompanied the poetry, but since they are found in both \Regius\ and \AM\ they must go back to a now-lost manuscript archetype.

The second layer is sts.~1--3 and 53--55, which together with the prose form a narrative frame dealing with King Garfrith and his imprisonment of Weden.

At the heart of the poem stand the gnomic wisdom stanzas 4--52, which make up the bulk of the poem.  The wisdom is mythological and sometimes obscure, and aligns with other Eddic verse like \textlink{Havamal}, \textlink{Vafthrudnismal}, \textlink{Baldrsdraumar}, \textlink{Sigrdrifumal}, and \textlink{Alvissmal}.

\subsection{Origin and motifs}

\Grimnismal\ consists at its core of the long speech given by the god Weden, and in this way it resembles \textlink{Havamal}.  Both poems may be said to fall squarely under the category of gnomic poetry, but whereas most of \textlink{Havamal} deals with life-advice of a more general type, \Grimnismal\ focuses entirely on mythological and cosmological lore.

The lore found in \Grimnismal\ is very varied and most of it will be discussed at length in the notes to the individual stanzas, but a few broader motifs ought to be discussed here.

\subsubsection{The wandering god}

Like in \textlink{Havamal} we here see Weden in his role as the god of travellers and strangers, for he himself is a stranger and frequently travels among mankind.  In \Grimnismal\ this is made explicit as the god sets out to test the hospitality of his protegé and punishes him harshly when he fails to uphold the duties of a king, which include generosity and hospitality.

The motif of a travelling god testing the hospitality of mankind has ancient and likely Indo-European origins.  It is in any case very prevalent in the related Greek tradition.  Already in Homer we find it said plainly, \Odyssey\ 17.483--487: “Antinoös, you did badly to hit the unhappy vagabond: / a curse on you, if he turns out to be some god from heaven. / For the gods do take on all sorts of transformations, appearing as / strangers from elsewhere, and thus they range at large through the cities, / watching to see which men keep the laws, and which are violent.” \parencite{LattimoreOd}.  Is it not just this wrath which we see in \Grimnismal, when in st.~55 the god, having assumed his true, terrible form, exclaims: “Draw near Me, if thou mayst!”

If any wandering stranger may be a god, indeed the highest god, it naturally follows that one should furnish the impoverished well, lest one incur the wrath of the god(s).  This in turn confers a level of heavenly protection on the hands of the traveller and makes hospitality a sacred duty.  As much is seen in the Greek epithet of Zeus, \textgreek{Ζεύς ξένιος} ‘Zeus of hospitality’ (Plato, \emph{Laws} 12.953e), the term used for the high god in His role as the protector of strangers and sacred hospitality (\textgreek{ξενῐ́ᾱ}).  With this in mind, we may rightly call Weden in this aspect *\emph{Wódanaz} \textgreek{ξένιος}.

\subsubsection{The land of the Gods}

TODO: the God-world, Shamanism, Ugdrassle’s Ash as the \emph{axis mundi}.  Cite: Eliade on shamanism, PCRN on the recreation of Osyard at a smaller scale.

\subsection{Summary}

The text begins with the frame narrative, which tells the story of the two young king’s sons Ayner and Garfrith.  The elder Ayner is fostered by Frie while the younger Garfrith is fostered by her husband, Weden.  After their father’s death, Garfrith betrays his elder brother upon the advice of Weden and becomes king after his father (P1).

One day Weden and Frie are arguing over their respective foster-sons and Frie accuses Garfrith of torturing wayfaring guests.  Weden sets out to test the hospitality of his protégé.  Unbeknownst to him, Frie has already sent her handmaid in disguise to warn Garfrith about the coming of a dangerous wizard.  When Weden arrives, he is promptly captured and tortured between two fires in order that he will reveal his name.  Garfrith’s young son, Ayner (clearly named after his uncle), kindly approaches the bound god and offers him a horn of drink.  Weden, assuming the name of Grimner, drinks from it and then takes to speak (P2).

Weden begins by complaining about the fires which are now burning his cloak (1); he states that for eight nights not a soul has offered him any help save Ayner, Garfrith’s son, who will soon become king after his father (2).  As thanks for the drink he gives him good health, and will offer him holy knowledge (3).

Here the gnomic section begins as Weden lists the individual abodes of the Gods (4--17).  The named Gods are Thunder (4), Wolder and Free (5), Weden and Sey (7), Weden (8--10), Shede (11), Balder (12), Homedall (13), Frow (14), Foresitter (15), Nearth (16), and Wider (17).

%This section has been discussed in detail by \textcite{deVries1952} TODO! who considers it corrupt. Specifically, he sees the second half of v. 4 as a later insert, since it does not elaborate on the “holy land” mentioned in the first half. \textcite{Jackson1995} argues convincingly against this, showing how the first half serves as a generalized introduction to the list; the holy land is the dwelling-places of the gods.)

After this list come several sts.~relating to Weden and his hall, Walhall (18--23).  Mentioned is the preparation of food in Walhall (18), Weden’s wolves (19) and ravens (20), the river through which the dead have to wade (21) and the gate through which they have to pass (22), the count of doors in Walhall (23), the count of doors in Thunder’s hall Bilshirner (24), and the goat and stag which stand on the hall and gnaw on the branches of the tree Leered (25--26).
From the horns of the stag Oakthirner water-droplets fall into Wharyelmer, which is the origin of all rivers (26).

Next follows a list of mythic rivers (27--28), ending with the waters through which Thunder must wade on his way to Ugdrassle (29).  This leads to a list of the horses ridden by the other gods on their way to Ugdrassle (31) which is followed by a description of the roots of Ugdrassle (31) and its animals (32--36).

Next follows some miscellaneous lore: a list of Walkirries (37), the horses that pull the sun (38), the shield that protects the earth from the sun’s heat (39) and the wolves that chase the sun and moon, the things created from Yimer’s body (41--42) with a digression on the significance of the \inx[P]{bloot} (43), the creation of the ship Shidebladner (44) and finally a list of the noblest of several categories of things and groups (45).

After these lists Weden utters an unclear stanza invoking the gods (46), before listing his many names (47--50).  He then turns to Garfrith, disappointed by the inhospitality and poor conduct of his former protégé, and foresees his imminent death (51--53).  He finally reveals himself by his true name, daring Garfrith to face him (53).  After this he repeats several of his names (54), and the poem ends.

Upon learning that the wizard is Weden himself, King Garfrith jumps up from his throne to take the god away from the fires, but trips and falls upon his unsheathed sword and dies.  After this his son Ayner becomes king and is said to have ruled the kingdom for a long time (P3).

\newpage

\section{Text}

\subsection{Frá sonum Hrauðungs konungs (‘From the sons of king Reading’)}

\bpg\bpa\mssnote{\Regius~8v/31, \AM~3v/23}%
Hrauðungr konungr átti tvá sonu. Hét annarr Agnarr, \edtext{en}{\Afootnote{add.~\emph{sonr} \AM}} \edtext{annarr}{\Afootnote{corr.~above line in same hand \AM}} Geir·røðr.
Agnarr var tíu vetra en Geir·røðr átta vetra. Þeir reru tveir á báti með dorgar sínar at smá-fiski.
Vindr rak þá í haf út. \edtext{Í nǫ́tt-myrkri brutu þeir}{\Afootnote{\emph{Þeir brutu í nǫ́tt-myrkri} \AM}} við land ok gingu upp; fundu \edtrans{kot-bónda einn}{a lone cottage farmer}{\Bfootnote{The farmer and his wife were Weden and Frie as is clarified by the following prose.}}.
Þar vǫ́ru þeir um vetr’inn.  Kerling fostraði Agnar, \edtrans{en \edtext{karl}{\Afootnote{add.~\emph{fostraði} \AM}} Geir·røð \edtext{ok kenndi hǫ́num rǫ́ð.}{\Afootnote{so \AM \emph{;} om.~\Regius}}}{but the farmer Garfrith and taught him counsel}{\Bfootnote{The motif of Weden favouring the youngest brother is also found in \textlink{Rigsthula}; for Weden’s counsel cf.~\textlink{Havamal}.}}.
At vári fekk karl þeim skip.  En er þau kerling leiddu þá til \edtext{strandar}{\Afootnote{\emph{skips} \AM.}}, \edtrans{þá mę́lti karl ein-mę́li við Geir·røð}{the husband spoke a private speech with Garfrith}{\Bfootnote{Instructing him to push his brother out to sea.}}.
Þeir fengu byr ok kvǫ́mu til stǫðva fǫður síns. Geir·røðr var fram í skipi.
Hann hljóp upp á land en hratt út skipi’nu, ok mę́lti: „\edtrans{Far þú \edtext{nú}{\Afootnote{so \AM \emph{;} om.~\Regius}} þar er smyl hafi þik!}{Go thou whither the fiends may have thee!}{\Bfootnote{The same curse is attested with \emph{gramir} for \emph{smyl} in \textlink{Harbardsljod}[60] and in a related form in \textlink{Brot}[10].}}“
Skip’it \edtext{rak}{\Afootnote{add.~\emph{í haf} \AM}} út en Geir·røðr gekk út til bǿjar. Hǫ́num var vel fagnat; \edtext{þá var faðir hans}{\Afootnote{\emph{en faðir hans var þá} \AM}} andaðr.
\edtext{Var þá Geir·røðr}{\Afootnote{\emph{Geir·røðr var þá} \AM}} til konungs tekinn, ok varð maðr á-gę́tr.\epa

\bpb {\huge K}ing Reading had two sons. One was called Ayner and the other Garfrith.
Ayner was ten winters old, but Garfrith eight winters. The two were rowing in a boat with their trolling-lines for small fishing.
The wind drove them out into the sea. In the dark of night they crashed onto land and walked ashore; they found a lone cottage farmer.
There they stayed over the winter. The farmer’s wife fostered Ayner but the farmer [fostered] Garfrith and taught him counsel.
In the spring the farmer gave them ships, but when he and his wife led them to the shore, the husband spoke private speech with Garfrith.
They caught good wind, and came to their father’s harbour. Garfrith was in the front of the ship.
He leapt onto land and pushed out the ship, and spoke: ”Go now whither the fiends may have thee!”
The ship drove out but Garfrith walked towards the farm.  He was welcomed well; by then was his father ended.
Garfrith was then taken as king, and became an excellent man.\epb\epg


\bpg\bpa\mssnote{\Regius~9r/10, \AM~4r/3}%
\Ballnote{The motif of Weden and Frie sitting in their Heavenly Seat, looking out over the world and each picking a different side in a conflict is also found in the short 7th century text the Origin of the Nation of the Longbeards (i.e., the Lombards; Latin \emph{Origo Gentis Langbardorum}).  The narrative is clearly based on oral tradition brought from Northern Europe in the first centuries of the current era and as such attests the great age of the present motif.  The relevant passage is edited at the end of the present poem under the Appendix.}%
Óðinn ok Frigg sǫ́tu í \edtrans{\edtext{Hlið-skjǫlfu}{\Afootnote{\emph{Hlið-skjǫlf} \AM}}}{the Lithshelf}{\Bfootnote{The lookout post of the Gods, which also appears in \textlink{Skirnismal}[P1].  Cf.~\Gylfaginning~9:11: \emph{Þar er einn staðr, er Hlið-skjǫlf heitir, ok þá er Óðinn setti⸗sk þar í hǫ́-sę́ti, þá sá hann of alla heima ok hvers manns at-hǿfi ok vissi alla hluti, þá er hann sá.} ‘Thereat (viz.~in Osyard) there is one place which is called the Lithshelf, and when Weden sat there in the high seat then he saw throughout all the Homes and each man’s behaviour, and knew all the things which he saw.’ and \Gylfaginning~17:8--9, cited in st.~6 n.~below.}} ok sǫ́ um \edtext{heima alla}{\Afootnote{rev.~\AM}}.
Óðinn mę́lti: „\edtrans{Sér þú Agnar fóstra þinn, hvar hann elr bǫrn við gýgi í helli’num}{Dost thou see Ayner, thy foster-son, where he begets children with a troll-woman in her cave?}{\Bfootnote{This may relate to Frie’s role as love-goddess.  Ayner is in any case a weak, effeminate man who devotes his life to copulation rather than war.}}?
\edtext{En Geir·røðr, fóstri minn, er konungr}{\Afootnote{\emph{En er konung fóstri minn} (corrupt) \AM}} ok sitr \edtext{nú}{\Afootnote{om. \AM}} at \edtext{landi}{\Afootnote{\emph{lǫndum} \AM}}.“
Frigg segir: „Hann er mat-níðingr sá at hann kvelr gesti sína ef hǫ́num þykkja of·margir koma.“
Óðinn segir at þat er in mesta lygi. Þau veðja um þetta mál.
Frigg sendi eskis-mey sína, Fullu, til Geir·røðar. Hón bað konung vara⸗sk at eigi fyr·gerði hǫ́num fjǫl-kunnigr maðr sá er þar var kominn í land, ok sagði þat mark á at engi hundr var svá ólmr at á hann myndi hlaupa.
En þat var inn mesti hé-gómi at Geir·røðr vę́ri eigi mat-góðr \edtext{ok}{\Afootnote{\emph{en} \AM}} þó lę́tr hann hand-taka þann mann er eigi vildu hundar á \edtext{ráða}{\Afootnote{\emph{hlaupa} \AM}}.
Sá var í feldi blǫ́m ok nefndi⸗sk Grímnir ok sagði ekki fleira frá sér þó’tt hann vę́ri at spurðr.
Konungr lét hann pína til sagna ok setja milli elda \edtext{tveggja}{\Afootnote{om. \AM}} ok sat \edtext{hann}{\Afootnote{om. \AM}} þar átta nę́tr.
Geir·røðr konungr \edtext{átti}{\Afootnote{add.~\emph{þá} \AM}} son tíu vetra gamlan ok hét Agnarr eptir bróður hans.
Agnarr gekk at Grímni ok gaf hǫ́num horn fullt at \edtext{drekka,}{\Afootnote{add.~\emph{ok} \AM}} sagði at \edtext{konungr}{\Afootnote{\emph{faðir hans} \AM}} gerði \edtext{illa}{\Afootnote{om.~\AM}} er hann \edtext{lét pína hann}{\Afootnote{\emph{píndi þenna mann} \AM}} sak-lausan.
Grímnir drakk af. Þá var eldr’inn svá kominn at feldr’inn brann af Grímni. Hann kvað:\epa

\bpb Weden and Frie sat in the \inx[P]{Lithshelf} and looked about all the Homes.
Weden spoke: “Dost thou see Ayner, thy foster-son, where he begets children with a troll-woman in her cave?
But Garfrith, \emph{my} foster-son, is king and now rules his land.”
Frie says: “He is such a meat-nithing that he torments his guests if he thinks too many are coming!”
Weden says that this is the greatest lie; they make a wager over this matter.
Frie sent her handmaid, Full, to Garfrith’s hall. She bade the king be wary, lest he be destroyed by the \inx[C]{manicunning} man who had come to his land; and said that his mark was that no hound was so fierce that it would rush at him.
But it was the greatest falsehood that Garfrith was not \inx[C]{good of meat}; and yet he has that man bound whom the hounds would not touch.
He was in a blue cloak and called himself Grimner, and did not tell anything more about himself, even though he was asked.
The king had him tortured that he would speak, and set him between two fires; and he sat there for eight nights.
King Garfrith had a son ten winters old, and he was called Ayner after his brother.
Ayner went up to Grimner and gave him a full horn to drink, saying that the king did badly as he had him tortured without cause.
Grimner drank it up. Then the fire had grown so much that the cloak burned on Grimner. Quoth he:\epb\epg\stepcounter{prosea}

\newpage

\subsection{The Speeches of Grimner}

\bvg\bva\mssnote{\Regius~9r/27, \AM~4r/17}%
\llap{„}\alst{H}ęitr est \alst{h}ripuðr \hld\ ok \alst{h}ęldr til mikill, &
\ind \edtrans{gǫng⸗umk \alst{f}irr}{go far from me}{\Bfootnote{Equivalent to \emph{gakk mér firr}.  The word \emph{gǫng⸗umk} ‘go \dots\ from me’ merits some linguistic discussion on the form of the imperative (\emph{gang-} rather than \emph{gakk}) and the suffixed pronoun \emph{⸗umk}.
The normal 2sg imper. of ON \emph{ganga} ‘to go’ is \emph{gakk}, which goes back to a sound change in PN whereby word-final homorganic clusters were devoiced, and in the case of clusters of the form \textipa{/\*N\*C/} assimilated into \textipa{/\r*{\*C}\textlengthmark/}.
Thus PN *\emph{gang} ‘go!’ (cf.~Got. \emph{gagg} \textipa{/gaNg/}, OHG \emph{gang}) > *\emph{gank} > \emph{gakk}, also *\emph{stand} ‘stand!’ > \emph{statt}, *\emph{bind} ‘bind!’ > \emph{bitt}.  This sound change was only active in early PN, and only affected those clusters which were word-final before syncope of unstressed vowels; thus in the ON 1sg pres. ind. \emph{ek gęng, ek stęnd} (< *\emph{ek gangu, ek standu}) ‘I go, I stand’ et c. we find the original cluster preserved.

At some time in early PN, accusative personal pronouns were suffixed to certain verbs; those which survive into ON are 1sg \emph{⸗umk} (= ON \emph{mik}) and 3sg reflexive \emph{⸗sk} (= ON \emph{sik}).  In later ON these clitics became the passive endings still found in modern North Germanic, but in archaic or poetic ON they could still serve as suffixed pronouns with a reflexive or even dative function, as is the case in this instance.  The fact that the suffixion of the pronoun has prevented assimilation of the consonant cluster \emph{ng} establishes its \emph{terminus ante quem} as the first syncope period, the early 7th century, and since it is unlikely that such an irregular form as \emph{gǫng⸗umk} could have survived for long alongside \emph{gakk}, its presence here probably allows us to conclude that \textlink{Grimnismal} is a rather old poem.

Other instances of dative \emph{⸗umk} include \textlink{Lokasenna}[35]/1a: \emph{es⸗umk}, \textlink{Fafnismal}[1]/4: \emph{stǫnd⸗umk}, \textlink{EddicFragments}[F2][II] 1/1: \emph{eru⸗mk}, and \Ragnarsdrapa\ 7/2a: \emph{gǫf⸗umk}.}} \alst{f}uni! &
\alst{L}oði sviðnar, \hld\ þó’tt ȧ \alst{l}opt \edtext{bera’k}{\Afootnote{\emph{vera} \AM}}; &
\ind brinn⸗umk \alst{f}eldr \alst{f}yrir.\eva

\bvb “{\huge H}\textsc{ot art thou, scorcher}, and far too great--- \\
\ind go far from me, O flame! \\
My wool cape is singed though I hold it aloft; \\
\ind the cloak burns before me!\evb\evg


\bvg\bva\mssnote{\Regius~9r/29, \AM~4r/18}%
\alst{Á}tta nę́tr \hld\ sat’k milli \alst{ę}lda hér, &
\ind svá’t mér \alst{m}ann⸗gi \alst{m}at né bauð &
nema \alst{ęi}nn Agnarr, \hld\ es \alst{ęi}nn skal ráða, &
\alst{G}ęir·røðar sonr, \hld\ \alst{G}otna landi.\eva

\bvb For eight nights I sat between the fires here, \\
\ind while no man offered me food, \\
save for Ayner alone, who alone shall rule--- \\
Garfrith’s son---the land of the Gots!\evb\evg


\bvg\bva\mssnote{\Regius~9r/31, \AM~4r/20}%
\alst{H}ęill skalt, Agnarr, \hld\ alls \alst{h}ęilan biðr &
\ind þik \alst{V}era-týr \alst{v}esa; &
\edtrans{\alst{ęi}ns drykkjar}{For a single drink}{\Bfootnote{The horn which Ayner offered Weden (P2 above).}} \hld\ skalt \alst{a}ldri⸗gi &
\ind \edtrans{bętri \alst{g}jǫld}{better recompense}{\Bfootnote{Namely the mythic lore which takes up sts.~4--45.}} \alst{g}eta:\eva

\bvb Hale shalt thou be, Ayner, for hale \\
\ind does Were-Tew \name{= Weden} bid thee be. \\
For a single drink shalt thou never get \\
\ind better recompense.\evb\evg


\bvg\bva\mssnote{\Regius~9r/33, \AM~4r/22}%
\alst{L}and es hęilagt, \hld\ es \alst{l}iggja sé’k &
\ind \alst{ǫ̇}sum ok \alst{ǫ}lfum nę́r; &
ęn ï \alst{Þ}rúð-hęimi \hld\ skal \alst{Þ}ȯrr vesa &
\ind \edtrans{und’s of \alst{r}júfa⸗sk \alst{r}ęgin}{until the Reins are ripped}{\Bfootnote{Until the \inx[P]{Rakes of the Reins}.  A formulaic expression; see note to \textlink{Baldrsdraumar}[14]/4 for further occurrences.}}.\eva

\bvb The land is holy which I see lying \\
\ind near the \inx[F]{Eese and Elves}, \\
but in Thrithham shall Thunder dwell \\
\ind until the Reins are ripped.\evb\evg


\bvg\bva\mssnote{\Regius~9v/2, \AM~4r/23}%
\alst{Ý}-dalir hęita, \hld\ þar’s \alst{U}llr hęfir &
\ind \alst{s}ér of gǫrva \alst{s}ali; &
\alst{A}lf-hęim Fręy \hld\ gǫ́fu ï \alst{á}r-daga &
\ind \alst{t}ívar at \edtrans{\alst{t}ann-féi}{tooth-gift}{\Bfootnote{The gift the child receives when he sheds his first tooth.}}.\eva

\bvb Yewdales they are called where Wolder has \\
\ind made for himself a hall. \\
Elfham to Free in days of yore \\
\ind the Tews as a tooth-gift gave.\evb\evg


\bvg\bva\mssnote{\Regius~9v/3, \AM~4r/25}%
\Ballnote{This stanza is paraphrased in \Gylfaginning~17:8--9: \emph{Þar er enn mikill staðr, er Vala-skjǫlf heitir.  Þann stað á Óðinn.  Þann gerðu goð’in ok þǫkðu skíru silfri, ok þar er Hlið-skjǫlf’in í þessum sal, þat há-sę́ti, er svá heitir, ok þá er Al-fǫðr sitr í því sę́ti, þá sér hann of alla heima.} ‘There is yet another great place which is called Waleshelf.  That place Weden owns.  The Gods made it and thatched it with pure silver, and in these hall is the Lithshelf, the high-seat which is called so, and when Allfather sits in that seat then he sees throughout all the Homes.’}%
\alst{B}ǿr es sá \edtrans{(hinn þriði)}{(the third)}{\Bfootnote{The first numbering of a location; the numbering continues up to st.~16.  Although parallelled by other numbered poetic lists (\textlink{Havamal}[147]--165, \textlink{Vafthrudnismal}[20]--42) some circumstances speak against the numbers in \Grimnismal\ being authorial.  First, the alliteration is never reliant on the numbers, which is quite different from the numbered items in \textlink{Havamal}[147]--165 and \textlink{Vafthrudnismal}[20]--42.  Secondly the numbering is inconsistent; Thunder’s realm (st.~4) is not counted, and Wider’s land (st.~17) has no numeral (perhaps since the form of the stanza would not allow it.)  Third, in \Gylfaginning’s citations of sts.~11--15 the numbers are missing.}}, \hld\ es \edtrans{\alst{b}líð ręgin}{the blithe Reins}{\Bfootnote{A reverent epithet for the Gods, also occurring in sts.~38/3, 42/1 below.  Parodied by Lock in \textlink{Lokasenna}[32]/3.  For other epithets cf.~\textlink{Vafthrudnismal}[13]/4 n.}} &
\ind \alst{s}ilfri þǫkðu \alst{s}ali; &
\alst{V}ala-skjǫlf hęitir, \hld\ \edtrans{es \alst{v}élti sér}{won through wiles}{\Bfootnote{Several previous editors and translators (e.g. \textcite{FinnurEdda}, \textcite{PettitEdda}, \textcite{LarringtonEdda}) have rendered this phrase with variants of “craftily made for himself”, where the verb \emph{véla} would mean ‘craftily make’.  To my knowledge this sense is never otherwise attested, and its common meaning is ‘defraud, trick, betray’.  A simpler reading would be to see this as a reference to the myth of the ettin Master Builder who built the wall of Osyard as alluded to in \textlink{Voluspa}[24]--25 and told at length in \Gylfaginning~42.  The Gods had promised him Frow and the Sun and Moon if he could build it within a year, but to slow him down they made use of various tricks.  When at last it seemed like he would make it in time, they called upon Thunder who arrived and killed him.}} &
\ind \alst{ǫ̇}ss ï \alst{á}r-daga.\eva

\bvb By is (the third) one, where the blithe Reins \\
\ind with silver thatched the halls. \\
Waleshelf is the name of the one which he won through wiles, \\
\ind the Os in days of yore.\evb\evg


\bvg\bva\mssnote{\Regius~9v/5, \AM~4r/26}%
\Ballnote{This stanza is likely paraphrased in \Gylfaginning~35:3: \emph{Ǫnnur er Sága.  Hón býr á Søkkva-bekk, ok er þat mikill staðr.} ‘The other [goddess] is Sey.  She lives in Sinkbench, and that is a great place.’}%
\alst{S}økkva-bękkr hęitir (hinn fjórði), \hld\ ęn þar \alst{s}valar knegu &
\ind \alst{u}nnir \edtext{\emph{glymja \alst{y}fir}}{\Afootnote{metr.~emend.; rev.~\Regius\AM}}; &
þar þau \alst{Ó}ðinn ok Sága \hld\ drekka umb \alst{a}lla daga &
\ind \alst{g}lǫð ór \alst{g}ollnum kęrum.\eva

\bvb Sinkbench is (the fourth) one called, but there do cool \\
\ind waves clash overhead; \\
there Weden and Sey drink for all days, \\
\ind glad, out of golden casks.\evb\evg


\bvg\bva\mssnote{\Regius~9v/7, \AM~4r/28}%
\Ballnote{For Gladsham see \Gylfaginning~14:2--7, which likely paraphrases this stanza, cited in \textlink{Voluspa}[7] n.}%
\alst{G}laðs-hęimr \edtext{hęitir}{\Afootnote{\emph{es} \AM}} (hinn fimti) \hld\ þar’s hin \alst{g}oll-bjarta &
\ind \alst{V}al-hǫll \alst{v}íð \edtext{of}{\Afootnote{om.~\AM}} þrumir; &
\edtext{ęn þar \alst{H}roptr \hld\ kýss \alst{h}vęrjan dag &
\ind \alst{v}ápn-dauða \alst{v}era.}{\lemma{ęn \dots\ vera ‘but \dots\ warriors.’}\Bfootnote{For Weden’s chosing of the slain see st.~14.}}\eva

\bvb Gladsham is (the fifth) one called, where the gold-bright \\
\ind Walhall, wide, stands fast, \\
but there Roft \name{= Weden} chooses every day \\
\ind weapon-dead warriors.\evb\evg


\bvg\bva\mssnote{\Regius~9v/10, \AM~4r/30}%
\Ballnote{In \Regius\ the order of sts.~9 and 10 is reversed, but a sign at the start of each suggests they should switch places; this is further indicated by the abbreviation of ll.~1--2 in st.~10.}%TODO: check mss again with pictures
Mjǫk ’s \alst{au}ð-kęnnt \hld\ þęim’s \edtext{til}{\Afootnote{om.~\AM}} \alst{Ó}ðins koma &
\ind \alst{s}al-kynni at \alst{s}éa, &
\edtrans{\alst{sk}ǫptum}{shafts}{\Bfootnote{Spear-shafts.}} ’s rann rępt, \hld\ \edtrans{\alst{sk}jǫldum ’s salr þakiðr}{with shields is the hall thatched}{\Bfootnote{The same is said in \Gylfaginning~2:6--7: \emph{Þak hennar var lagt gylldum skjǫldum svá sem spán-þak.  Svá segir Þjóð·ólfr inn hvin-verski, at Val-hǫll var skjǫldum þǫkð.} ‘Its roof was covered with gilt shields like a shingle-roof.  So does Thedwolf the Whine-dweller say that Walhall was thatched with shields.’  After this, \Gylfaginning\ cites Þhorn \emph{Harkv} 11/1--2, which uses the kenning \emph{Sváfnis sal-nę́frar} ‘Swebner’s hall-shingles \ken{shields}’, where Swebner is a name for Weden (cf.~st.~57 below).}}, &
\ind \alst{b}rynjum of \alst{b}ękki stráat.\eva

\bvb Very easily recognized by those who come to Weden \\
\ind is the hall to see: \\
With shafts is the house roofed, with shields is the hall thatched; \\
\ind with byrnies the benches strewn.\evb\evg


\bvg\bva\mssnote{\Regius~9v/9, \AM~4r/31}%
Mjǫk ’s \alst{au}ð-kęnnt \hld\ þęim’s til \alst{Ó}ðins koma &
\ind \edtext{\alst{s}al-kynni at \alst{s}éa}{\Afootnote{\emph{†sia at sia†} \AM}}, &
\edtext{\alst{v}argr hangir \hld\ fyr \alst{v}estan dyrr &
\ind ok drúpir \alst{ǫ}rn \alst{y}fir.}{\lemma{vargr hangir \hld\ fyr vestan dyrr / ok drúpir ǫrn yfir. ‘A wolf hangs before the western door, and an eagle droops down over it.’}\Bfootnote{Something very similar is found in Widukind of Corvey’s \emph{History of the Saxons} 1:12.  The Saxons had just conquered a fortress, and \emph{mane [...] facto ad orientalem portam ponunt aquilam, aramque victoriae construentes secundum errorem paternum sacra sua propria veneratione venerati sunt} ‘at the coming of morning they set an eagle at the eastern gate, and, building an altar of victory, they worshipped it with their own holy worship in accordance with their ancestral error.’  The altar was pledged to \inx[P]{Ermen}, whom Widukind identifies with Mars or Hermes, certainly Weden.  According to \textcite[156]{HyltenCavallius1863} it was custom in Wärend, southern Sweden to hang the bodies of killed wolves high up in old oaks, and killed birds of prey above the stable-door.}}\eva

\bvb Very easily recognized by those who come to Weden \\
\ind is the hall to see: \\
A wolf hangs before the western door, \\
\ind and an eagle droops down over it.\evb\evg%TODO: bibliography for Widukind


\bvg\bva\mssnote{\Regius~9v/12, \AM~4v/2, \\
\RegiusProse~6v/20, \Trajectinus~TODO, \\
\Wormianus~TODO, \Upsaliensis~TODO}%
\Ballnote{St.~11 is cited in \Gylfaginning~23:15, introduced by the words \emph{Svá er sagt} ‘So it is said’ and concluding the chapter.  For the rest of the chapter see \textlink{EddicFragments}[F2][II].}%
\edtext{\alst{Þ}rym-hęimr}{\Afootnote{\emph{Þrúð-hęimr} \Upsaliensis}} hęitir \edtrans{(hinn sétti)}{the sixth}{\Afootnote{om. \GylfMS}}, \hld\ \edtrans{es}{where}{\Afootnote{\emph{þar nú} ‘where now’ \Upsaliensis}} \alst{Þ}jatsi \edtrans{bjó}{dwelled}{\Afootnote{om. \Wormianus; \emph{býr} ‘dwells’ \Upsaliensis}}, &
\ind sá hinn \edtrans{\edtext{\alst{ȧ}m-átki}{\Afootnote{\emph{mátki} \Upsaliensis}} \alst{jǫ}tunn}{uncanny ettin}{\Bfootnote{Formulaic. See note to \textlink{Voluspa}[8].}}; &
ęn nú \alst{Sk}aði byggvir, \hld\ \alst{sk}ír brúðr \edtrans{goða}{of the Gods}{\Afootnote{\emph{guma} ‘of men’ \Upsaliensis}}, &
\ind \edtext{\alst{f}ornar}{\Afootnote{\emph{forna} \AM \emph{;} \emph{†førnir†} \Trajectinus}} toptir \alst{f}ǫður.\eva

\bvb Thrimham is (the sixth) one called, where Thedse dwelled, \\
\ind that uncanny ettin; \\
but now Shede bedwells---the pure bride of the Gods--- \\
\ind the ancient plots of her father.\evb\evg


\bvg\bva\mssnote{\Regius~9v/14, \AM~4v/3, \\
\RegiusProse~6v/4, \Trajectinus~TODO, \\
\Wormianus~TODO, \Upsaliensis~TODO}%
\Ballnote{St.~12 is cited in \Gylfaginning~22:22, introduced by the words \emph{Svá sem hér segir} ‘Just as it says here’ and concluding the chapter, which tells about Balder.  \Gylfaginning~22:21 provides context: \emph{Hann býr þar, sem heitir Breiða-blik.  Þat er á himni.  Í þeim stað má ekki vera ó·hreint.} ‘He [viz.~Balder] lives at the place called Broadblicks.  That is in Heaven.  In that place may nothing be unclean.’}%
\edtext{\alst{B}ręiða-\alst{b}lik}{\Afootnote{\emph{Bręiða} (with \emph{blik} added above the line in other hand) \Regius \emph{;} \emph{Bręiða-†blio†} \AM}} \edtrans{eru (hin sjaundu)}{are (the seventh)}{\Afootnote{\emph{hęita} ‘[they] are called’ \RegiusProse\Trajectinus\Wormianus \emph{;} \emph{hęitir} ‘is one called’ \Upsaliensis.}}, \hld\ ęn þar \alst{B}aldr hęfir &
\ind \alst{s}ér of gǫrva \alst{s}ali, &
ȧ því \alst{l}andi \hld\ es \alst{l}iggja vęit’k &
\ind \alst{f}ę́sta \edtrans{\edtext{\alst{f}ęikn-stafi}{\Afootnote{\emph{†fæing†-stafi} \Upsaliensis}}}{wicked deeds}{\Bfootnote{This sense of \emph{stafir} (lit.~‘staffs, staves’) is common in poetry.  Cf.~the direct cognate in \Beowulf\ 1018b: \emph{fâcen-stafas}, generally taken as referring to treacherous intrigues among the \inx[P]{Shieldings} \parencite[177]{KlaeberBeowulf}.}}.\eva

\bvb Broadblicks are (the seventh), but there Balder has \\
\ind made for himself a hall, \\
in that land where I know lie \\
\ind the fewest wicked deeds.\evb\evg


\bvg\bva\mssnote{\Regius~9v/16, \AM~4v/5, \\
\RegiusProse~7r/21, \Trajectinus~TODO, \\
\Wormianus~TODO, \Upsaliensis~TODO}%
\Ballnote{St.~13 is cited in \Gylfaginning~27:14, the chapter dealing with Homedall, introduced by the words: \emph{Hér er svá sagt} ‘Here it is said so’.  For the full chapter see \textlink{EddicFragments}[F3][III].}%
\alst{H}imin-bjǫrg \edtrans{eru (hin ǫ́ttu)}{are (the eighth)}{\Afootnote{\emph{hęita} ‘[they] are called’ \RegiusProse\Trajectinus\Wormianus \emph{;} \emph{hęitir} ‘is one called’ \Upsaliensis.}}, \hld\ ęn þar \edtext{\alst{H}ęim·dall}{\Afootnote{\emph{Hęim·dallr býr} \Upsaliensis}} &
\ind kveða \alst{v}alda \alst{v}éum; &
\edtext{þar}{\Afootnote{om.~\Upsaliensis}} \edtrans{\edtext{\alst{v}ǫrðr}{\Afootnote{\emph{vǫrðum} \Upsaliensis}} goða}{Watchman of the Gods}{\Bfootnote{Formulaic epithet of Homedall, also occurring in \textlink{Lokasenna}[48]/4 and possibly in \textlink{Skirnismal}[28] (\emph{vǫrðr með goðum} ‘the Watchman among the Gods’).  A similar epithet is found in \Husdrapa\ 2, \emph{vári ragna} ‘Defender of the Reins’.  \Gylfaginning~27, where the present stanza is cited, gives some further details; see \textlink{EddicFragments}[F3][III].}} \hld\ drekkr ï \alst{v}ę́ru ranni &
\ind \alst{g}laðr \edtext{hinn}{\Afootnote{so \AM\GylfMS \emph{;} om. \Regius}} \alst{g}óða mjǫð.\eva

\bvb Heavenbarrows are (the eighth), but there Homedall, \\
\ind they say, wields over wighs. \\
There the Watchman of the Gods drinks in the tranquil house, \\
\ind glad, the good mead.\evb\evg


\bvg\bva\mssnote{\Regius~9v/17, \AM~4v/6, \\
\RegiusProse~6v/27, \Trajectinus~TODO, \\
\Wormianus~TODO, \Upsaliensis~TODO}%
\Ballnote{St.~14 is cited in \Gylfaginning~24:9, introduced by the words \emph{svá sem hér segir} ‘just as it says here’. The immediately preceding \Gylfaginning~24:6--8 provides context: \emph{En Freyja er á·gę́tust af Ǫ́synjum.  Hón á þann bǿ á himni er Fólk-vangar heita, ok hvar sem hón ríðr til vígs þá á hón halfan val, en halfan Óðinn.} ‘But Frow is the most distinguished of the Ossens.  She has the dwelling in Heaven which is called the Folkwongs, and wherever she rides into battle she owns half of the slain, but Weden [has the other] half’.}%
\alst{F}olk-vangr \edtrans{es (hinn níundi)}{is (the ninth)}{\Afootnote{\emph{hęitir} ‘[one] is called’ \GylfMS}}, \hld\ ęn þar \alst{F}ręyja rę́ðr &
\ind \edtext{\alst{s}essa kostum ï}{\Afootnote{\emph{kosta bętstum} \Upsaliensis}} \alst{s}al; &
\edtrans{\alst{h}alfan val \hld\ hǫ̇n kýss \edtext{\alst{h}vęrjan dag}{\Afootnote{\emph{hvęrn dag} \Upsaliensis}}}{half the slain she chooses each day}{\Bfootnote{Apart from the paraphrase in \Gylfaginning~this statement is unprecedented; in all other sources the battle-slain are devoted exclusively to Weden and chosen by his walkirries.

The roots of \emph{kjósa val} ‘choose the slain’ are the same as those in \emph{val-kyrja} ‘walkirrie’, lit.~‘chooser of the slain’), and as Frow is a prominent goddess this could be seen to make her the chief walkirrie.  This is parallelled by \Sorlathattr, where Frow assumes the name \inx[C]{Gandle} (\emph{Gǫndul}, a name attested in several lists of walkirries; see \textlink{Voluspa}[30] n.) and incites the legendary never-ending Conflict of the Headenings (\emph{Hjaðningavíg}).

In spite of this there is reason to believe that the chief walkirrie was rather \inx[C]{Frie}, Weden’s wife.  First, one of the functions of the Walkirries is to bear ale to the Oneharriers (\textlink{Grimnismal}[37]).  This mirrors a real custom at Germanic banquets attested where the host’s wife or daughter would pour ale to his retainers and guests (the so-called “lady with a mead cup” ritual most famously seen in \Beowulf; see \textcite{Enright1996} and \textcite{Riseley2014}).  As Weden’s wife, we would expect Frie to have this role.
Second, at Balder’s funeral (\Gylfaginning~49:49--52), Frie rides with Weden and the Walkirries while Frow rides alone with her cats; if Frow were the chief walkirrie we would probably expect her to ride at the head of them.
Third, there are two separate myths where Frie and Weden contend over the fates of armies and men; these are the prose introduction to the present poem (P2 above) and the Longbeard origin myth (see Appendix below), whereas Frow and Weden do not otherwise come into conflict over the slain.

Whether the present stanza originally referred to Frie is hard to say, although confusion between Frow and Frie would probably not be expected in a source as old as \Grimnismal.}}, &
\ind ęn halfan \alst{Ó}ðinn \alst{á}.\eva

\bvb Folkwong is (the ninth), but there Frow rules \\
\ind the choice of seats in the hall; \\
half the slain she chooses each day, \\
\ind but half does Weden own.\evb\evg


\bvg\bva\mssnote{\Regius~9v/19, \AM~4v/8, \\
\RegiusProse~7r/34, \Trajectinus~TODO, \\
\Wormianus~TODO, \Upsaliensis~TODO}%
\Ballnote{St.~15 is cited in \Gylfaginning~32:4, introduced by the words \emph{Svá segir hér.} ‘So it says here.’  \Gylfaginning~32:2--3 provides context: \emph{Hann á þann sal á himni, er Glitnir heitir.  En allir, er til hans koma með sakar-vand-rę́ði, þá fara allir sáttir á braut.  Sá er dóm-staðr betstr með goðum ok mǫnnum.} ‘He has the hall in heaven which is called Glitner.  But all who come to him with difficult disputes then all go away reconciled.  That is the best judgment place amongst gods and men.’}%
\alst{G}litnir \edtrans{es (hinn tíundi)}{is (the tenth)}{\Afootnote{\emph{‘h. e. x.’} (for *\emph{‘e. h. x.’}) \AM \emph{;} \emph{hęitir salr} ‘a hall is called’ \GylfMS}}, \hld\ hann \edtext{’s}{\Afootnote{om.~\Upsaliensis}} \alst{g}olli studdr &
\ind ok \alst{s}ilfri \edtext{þakðr it}{\Afootnote{om.~\Upsaliensis}} \alst{s}ama; &
ęn \edtext{þar}{\Afootnote{\emph{þat} \RegiusProse\Trajectinus \emph{;} \emph{þá} \Wormianus}} \alst{F}or·seti \hld\ byggvir \alst{f}lęstan dag &
\ind ok \alst{s}vę́fir allar \alst{s}akir.\eva

\bvb Glitner is (the tenth): it is supported by gold, \\
\ind and thatched with silver likewise, \\
but there Foresitter dwells for most days, \\
\ind and puts all disputes to sleep.\evb\evg


\bvg\bva\mssnote{\Regius~9v/21, \AM~4v/9}%
\alst{N}óa-tu̇n eru (hin ęlliptu), \hld\ ęn þar \alst{N}jǫrðr hęfir &
\ind \alst{s}ér of gǫrva \alst{s}ali; &
\edtrans{\alst{m}anna þęngill \hld\ hinn \alst{m}ęins-vani}{The lord of men, the guiltless one}{\Bfootnote{Interesting epithets probably relating to Nearth’s roles in upholding the prosperity of the land and the law.  Cf.~my article on pre-Christian oaths (TODO).}} &
\ind \edtrans{\alst{h}ǫ́-timbruðum \alst{h}ǫrgi rę́ðr}{rules the harrow timbered on high}{\Bfootnote{The rare verb \emph{hǫ́-timbra} ‘timber on high’ otherwise only occurs in \textlink{Voluspa}[7], likewise in connection with the \emph{hǫrgr} ‘harrow’.  The harrow is an outdoors holy place; see there.  Cf.~also \textlink{Vafthrudnismal}[38] where Nearth is said to rule a great many \emph{hofum ok hǫrgum} ‘hoves and harrows’.}}.\eva

\bvb Nowetons are (the eleventh), and there Nearth has \\
\ind made for himself a hall. \\
The lord of men, the guiltless one, \\
\ind rules the \inx[C]{harrow} timbered on high.\evb\evg


\bvg\bva\mssnote{\Regius~9v/23, \AM~4v/11}%
\Ballnote{At the Rakes of the Reins Wider must avenge his father, Weden.  See \textlink{Voluspa}[51]--52, \textlink{Vafthrudnismal}[53].}%
\edtrans{\alst{H}rísi vęx \hld\ ok \alst{h}ǫ́u grasi}{With brushwood and with tall grass grows}{\Bfootnote{Identical to \textlink{Havamal}[119]/6.}} &
\ind \alst{V}íð·ars land, \alst{v}iði, &
ęn þar \alst{m}ǫgr of lę́t⸗sk \hld\ af \alst{m}ars baki &
\ind \alst{f}rǿkn \edtext{\emph{at}}{\Afootnote{\emph{ok} \emph{(‘⁊’)} \Regius\AM}} hęfna \alst{f}ǫður.\eva

\bvb With brushwood and with tall grass grows, \\
\ind \inx[P]{Wider}’s land, with wood, \\
but there the lad vows from the back of his steed, \\
\ind brave, to avenge his father.\evb\evg


\bvg\bva\mssnote{\Regius~9v/24, \AM~4v/12, \\
\RegiusProse~9v/12, \Trajectinus~10r/18, \\
\Wormianus~TODO, \Upsaliensis~TODO}%
\Ballnote{The poem moves away from the dwellings of the individual gods to an extended description of Walhall and its surroundings.  The cook Andrimner ‘Face-sooty’ cooks the boar Sowrimner ‘Sow-sooty’ in the cauldron Eldrimner ‘Fire-sooty’; by this meat are the Oneharriers nourished. ---

Sts.~18--20 are cited in order in \Gylfaginning~38 interspersed by short prose explanations.  St.~18 is cited following \Gylfaginning~38:6.  \Gylfaginning~38:3, 5--6 introduces it: \emph{En aldri er svá mikill mann-fjǫlði í Val-hǫll at eigi má þeim enda⸗sk flesk galtar þess er Sę́·hrímnir heitir.  Hann er soðinn hvern dag ok heill at aptni.  [...]  And·hrímnir heitir steikari’nn en Eld·hrímnir ketill’inn.  Svá er hér sagt:} ‘But there is never such a great gathering of men in Walhall that they should not be fully furnished with the flesh of that boar which is called Sowrimner.  He is cooked every day and is healthy and alive in the evening.  [...]  Andrimner is the name of the cook, but Eldrimner the kettle.  So it is said here:’}%
\edtext{\alst{A}nd·hrïmnir}{\Afootnote{\emph{And·runnir} \Trajectinus}} \hld\ lę́tr ï \edtext{\alst{Ę}ld·hrïmni}{\Afootnote{\emph{Ęld·runni} \Trajectinus}} &
\ind \edtext{\alst{S}ę́·hrïmni}{\Afootnote{\emph{Sę́·runni} \Trajectinus}} \alst{s}oðinn, &
\edtext{\alst{f}lęska}{\Afootnote{\emph{†felska†} \RegiusProse}} bętst, \hld\ ęn þat \alst{f}áir vitu, &
\ind \edtext{við}{\Afootnote{so \GylfMS \emph{;} om.~\Regius\AM}} hvat \alst{ę}in-hęrjar \alst{a}la⸗sk.\eva

\bvb Andrimner lets in Eldrimner \\
\ind Sowrimner be cooked, \\
the best of meats, but few know this: \\
\ind by what the \inx[P]{Oneharriers} are nourished.\evb\evg


\bvg\bva\mssnote{\Regius~9v/26, \AM~4v/14, \\
\RegiusProse~9v/17, \Trajectinus~10r/23, \\
\Wormianus~TODO, \Upsaliensis~TODO}%
\Ballnote{St.~19 is cited in \Gylfaginning~38:10, introduced by the words \emph{Svá segir hér.} ‘So it says here.’  \Gylfaginning~38:7--9 provides context: \emph{Þá mę́lti Gangleri: „Hvárt hefir Óðinn þat sama borð-hald sem Ein-herjar?“ Hárr segir: „Þá vist, er á hans borði stendr, gefr hann tveim úlfum, er hann á, er svá heita, Geri ok Freki.  En enga vist þarf hann: Vín er honum bę́ði drykkr ok matr.} ‘Then spoke Gangler: “Does Weden have the same things on his table as the Oneharriers?” High says: “The foodstuff which stands on his table he gives to the two wolves he owns, which are called thusly: Gar and Freak.  But he needs no foodstuffs; for him wine is both drink and food.’}%
\edtrans{\alst{G}era ok Freka}{Gare and Freak}{\Bfootnote{Weden’s two hounds.  It is not said with what food Weden feeds them, but if we consider Weden’s relation to the common ‘beasts of battle’ motif, it is most likely the corpses of dead warriors.}} \hld\ sęðr \edtext{\alst{g}unn-tamiðr}{\Afootnote{\emph{†gunnta midr†} \Trajectinus}}, &
\ind \alst{h}róðigr \alst{H}ęrja-fǫðr, &
ęn við \edtrans{\alst{v}ïn}{wine}{\Bfootnote{The wine on which Weden himself subsists may perhaps be identified with drink offerings.  For hallowing of alcohol to this god cf.~the 7th c. \emph{Vita Columbani}, ch. 53, describing a heathen rite of the Swabians (\cite[1040D--1041A]{VitaColumbani}, my translation):

\textlatin{Quo cum moraretur, et inter habitatores loci illius progrederetur, reperit eos sacrificium profanum litare velle, vasque magnum, quod vulgo cupam vocant, quod viginti et sex modios amplius minusve capiebat, cervisia plenum in medio habebant positum. Ad quod vir Dei accessit, et sciscitatur quid de illo fieri vellent. Illi aiunt deo suo Vodano, quem Mercurium vocant alii, se velle litare.}
‘While he was staying there and going about the dwellers of that place, he found out that they were going to offer a profane sacrifice, and a large cask called a \emph{cupa}, which held about twenty-six measures, was filled with beer and set in their midst.  When the man of God asked what they wanted to do with it they answered that they wanted to offer it to their god \emph{Vodan}, whom others call Mercury.’}} ęitt \hld\ \edtext{\alst{v}ápn-gǫfugr}{\Afootnote{\emph{vápn-gafigr} \RegiusProse\Trajectinus \emph{;} \emph{vápn-†gaffiþr†} \Upsaliensis}} &
\ind \alst{Ó}ðinn \edtext{\alst{ę́}}{\Afootnote{\emph{er} \Trajectinus}} lifir.\eva

\bvb \inx[P]{Gar and Freak} does the battle-accustomed \\
\ind glorious Father of Hosts \name{= Weden} feed, \\
but on wine alone, esteemed of weapons, \\
\ind Weden ever lives.\evb\evg


\bvg\bva\mssnote{\Regius~9v/28, \AM~4v/15, \\
\RegiusProse~9v/22, \Trajectinus~10r/28, \\
\Wormianus~TODO, \Upsaliensis~TODO}%
\Ballnote{St.~20 is cited in \Gylfaginning~38:13 and concludes the chapter, introduced by the words \emph{Svá sem sagt er} ‘Just as it is said’.  \Gylfaginning~38:11--12 provides context: \emph{Hrafnar tveir sitja á ǫxlum hónum ok segja í eyru hónum ǫll tíðendi, þau er þeir sjá eða heyra.  Þeir heita svá, Huginn ok Muninn.  Þá sendir hann í dagan at fljúga um heim allan, ok koma þeir aptr at dǫgurðar-máli.  Þar af verðr hann margra tíðenda víss.  Því kalla menn hann Hrafna-guð.} ‘Two ravens sit upon his shoulders and tell into his ears all the tidings which they see or hear.  They are called so: Highen and Minden.  He sends them out at dawn to fly throughout all the Home, and they come back at supper-time; thereof he is made aware of many tidings.  For this reason men call him Raven-god.’}%
\alst{H}uginn ok Muninn \hld\ fljúga \edtext{\alst{h}vęrjan dag}{\Afootnote{\emph{hvęrn dag} \Upsaliensis}} &
\ind \edtrans{\alst{jǫ}rmun-grund}{ermen-ground}{\Bfootnote{‘The world-ground’, i.e., ‘the whole wide world’, a poetic compound denoting the earth as a flat, spacious expanse of land, surrounded by the unbounded outer sea.  In ON this word only occurs elsewhere at two places: in a kenning on the late C10th Karlevi stone (Öl 1), which refers to the unbounded sea as \emph{Ęndils jǫrmun-grund} ‘Andle’s ermen-ground’ (Andle being a known sea-king), and in the C13th Sturl \emph{Hryn} 16/2 (\Skp\ 2), referring generically to land.  A direct cognate, OE \emph{eormen-grund}, appears only once, in \Beowulf\ 859, where it carries the same sense as in the present stanza.}} \alst{y}fir; &
\edtext{\alst{ó}⸗umk}{\Afootnote{\emph{otunk} \Trajectinus \emph{;} \emph{unz} \Upsaliensis}} \edtext{of}{\Afootnote{\emph{ek} \GylfMS}} Hugin, \hld\ %
at \alst{a}ptr \edtext{né}{\Afootnote{om.~\Trajectinus\Upsaliensis}} \edtext{komi⸗t}{\Afootnote{\emph{komi} \AM\RegiusProse\Trajectinus\Wormianus \emph{;}
\emph{kømr} \Upsaliensis}}; &
\ind þó sé⸗umk \alst{m}ęirr of \alst{M}unin.\eva

\bvb Highen and Minden fly every day \\
\ind over the ermen-ground \ken{earth}. \\
I worry for Highen that he might not come back, \\
\ind yet I fear more for Minden.\evb\evg


\bvg\bva\mssnote{\Regius~9v/30, \AM~4v/17}%
\Ballnote{A difficult stanza.  Considering the st.~is in the middle of others describing Weden and Walhall, Thound is possibly the river moat which the dead warriors have to pass over to reach Walhall.  The stanza may also be referring to the punishment of criminals in waters (for which see \textlink{Voluspa}[38] n.).}%
\alst{Þ}ýtr \alst{Þ}und, \hld\ unir \edtext{\alst{Þ}jóð·vitnis &
\ind \alst{f}iskr}{\lemma{Þjóð·vitnis fiskr ‘Thedwitner’s fish’}\Bfootnote{\emph{Þjóð·vitnir} is easily analyzed as \emph{þjóð-} ‘great, main’ + \emph{vitnir} ‘wolf’ and the great wolf is of course the \inx[P]{Fenrerswolf}.  The Wolf is called the brother of the Middenyardswyrm (\textlink{Hymiskvida}[23]), and so that may be its “fish”.  That the Wyrm can be called a fish is shown by \textlink{Hymiskvida}[24].  The “fish” can also be Nithehewer, the dragon who tortures slain criminals (\textlink{Voluspa}[38]).}} \alst{f}lóði ï; &
\alst{á}ar straumr \hld\ \edtext{þykkir}{\Afootnote{so \AM \emph{;} om. \Regius}} \alst{o}f-mikill &
\ind \edtext{\alst{v}al-glaumi}{\Afootnote{\emph{val-glaumni} \AM}} at \alst{v}aða.\eva

\bvb \inx[P]{Thound} roars; Thedwitner’s fish \\
\ind thrives in the flood. \\
The river-stream seems far too great \\
\ind for the noisy slain host to wade.\evb\evg


\bvg\bva\mssnote{\Regius~9v/32, \AM~4v/18}%
\edtrans{\alst{V}al-grind}{Walgrind}{\Bfootnote{‘Slain-gate’, the gate standing before Walhall.}} hęitir \hld\ es stęndr \alst{v}ęlli ȧ &
\ind \alst{h}ęilǫg fyr \alst{h}ęlgum durum; &
\alst{f}orn ’s sú grind, \hld\ ęn þat \alst{f}áir vitu, &
\ind hvé hǫ̇n ’s ï \alst{l}ȧs \edtext{of}{\Afootnote{so \AM \emph{;} om. \Regius}} \alst{l}okin.\eva

\bvb \inx[P]{Walgrind} it is called which stands on the plain, \\
\ind holy, before the holy doors. \\
Old is that gate, but few know this: \\
\ind how its lock is locked.\evb\evg


\bvg\bva\mssnote{\Regius~9v/34, \AM~4v/22, \\
\RegiusProse~6r/24, \Trajectinus~TODO, \\
\Wormianus~TODO, \Upsaliensis~TODO}%
\Aallnote{Sts.~23 and 24 come in reverse order in \AM.}%
\Ballnote{This st.~is cited in \Gylfaginning~21, the chapter describing Thunder, immediately following \Gylfaginning~21:4--6, which reads: \emph{Hann á þar ríki er Þrúð-vangar heita, en hǫll hans heitir Bil-skirnir.  Í þeim sal eru fimm hundrað gólfa ok fjórir tigir.  Þat er hús mest svá at menn viti.  Svá segir í Grímnis-mǫ́lum:} ‘He has his domain in the place called the Thrithwongs, but his hall is called Bilshirner.  In that hall there are five hundred and fourty floors.  That is the greatest house which men know of.  So it says in the Speeches of Grimner:’.}%
\alst{F}imm hundruð golfa \hld\ ok \edtext{umb}{\Afootnote{om.~\Upsaliensis}} \alst{f}jórum tøgum &
\ind svá \edtext{hygg’k}{\Afootnote{add.~\emph{á Val-hǫll vesa} (dittography through homoioarcton) \AM}} \alst{B}il-skirni með \alst{b}ugum; &
\alst{r}anna þęira, \hld\ es ek \edtext{\alst{r}ępt}{\Afootnote{\emph{rę́fr} \Upsaliensis}} vita, &
\ind \edtext{\alst{m}ïns}{\Afootnote{\emph{†murs†} \Trajectinus}} vęit’k męst \alst{m}agar.\eva

\bvb With five hundred floors and around fourty--- \\
\ind so I judge \inx[P]{Bilshirner} altogether. \\
Of those houses which I might know rafted \\
\ind I know my lad’s \ken*{= Thunder} is the greatest.\evb\evg


\bvg\bva\mssnote{\Regius~10r/2, \AM~4v/20, \\
\RegiusProse~10r/10, \Trajectinus~TODO, \\
\Wormianus~TODO, \Upsaliensis~TODO}%
\Ballnote{This st.~is cited in \Gylfaginning~40:4, introduced by the words: \emph{Hér mátt-u heyra í Grímnis-mǫ́lum.} ‘Here mightst thou hear in the Speeches of Grimner.’  \Gylfaginning~40:1--3 provides the context \emph{Þá mę́lti Gangleri: „Þetta eru undar⸗lig tíðendi, er nú sagðir þú.  Geysi-mikit hús mun Val-hǫll vera.  All-þrøngt mun þar opt vera fyrir durum.“ Þá svarar Hárr: „Hví spyrr þú eigi þess, hversu margar dyrr eru á hǫll’inni eða hversu stórar?  Ef þú heyrir þat sagt, þá munt⸗u segja, at hitt er undar⸗ligt, ef eigi má ganga út ok inn hverr, er vill.  En þat er með sǫnnu at segja, at eigi er þrøngra at skipa hana en ganga í hana.} ‘Then spoke Gangler: “These are wondrous tidings which thou hast now told.  An enormously great house must Walhall be.  It must often be fully cramped before the doors there.”  Then answers High: “Why askest thou not how many doors there are on the hall or how great they are?  If thou hearest that told then thou wilt say that the contrary is wondrous, that every one should not be able to go out and in as he wants to.  But with truth it is to be said that it is not more crowded to be inside it than to enter it.’}%
\alst{F}imm \edtext{hundruð}{\Afootnote{\emph{hund} \Trajectinus}} dura \hld\ ok \edtext{umb}{\Afootnote{om. \AM\Upsaliensis}} \edtext{\alst{f}jórum tøgum}{\Afootnote{abbrev.~as \emph{.xl.} \AM \emph{;} \emph{fjóra tugu}}}, &
\ind svá \edtext{hygg}{\Afootnote{\emph{kveð’k} \AM}} \edtext{ȧ}{\Afootnote{\emph{at} \Regius \emph{;} om.~\Trajectinus}} \edtext{\alst{V}al-hǫllu}{\Afootnote{\emph{Val-hǫll} \AM\Trajectinus\Wormianus}} \alst{v}esa; &
\edtrans{\alst{á}tta hundruð}{eight hundred}{\Bfootnote{The hundred is probably here the long hundred (120, rather than 100), which gives a sum of 640 * 960 = 614~400 Oneharriers.}} \alst{Ę}in-hęrja \hld\ ganga \edtext{sęnn}{\Afootnote{om.~\Regius\Upsaliensis}} ór \alst{ęi}num durum, &
\ind \edtrans{þȧ’s \edtext{fara}{\Afootnote{\emph{ganga} \Upsaliensis}} \edtext{við}{\Afootnote{so \AM\Wormianus\Upsaliensis \emph{;} \emph{at} \Regius \emph{;} \emph{með} \RegiusProse\Trajectinus}} \alst{v}itni at \alst{v}ega}{when to fight with the Wolf they fare}{\Bfootnote{At the Rakes of the Reins by Weden’s side.
Cf.~\textlink{Voluspa}[51]--52, \textlink{Vafthrudnismal}[53].}}.\eva

\bvb Five hundred doors and around fourty--- \\
\ind so I judge there to be on Walhall. \\
Eight hundred \inx[P]{Oneharriers} march at once out of one door \\
\ind when to fight with the Wolf they fare.\evb\evg


\bvg\bva\mssnote{\Regius~10r/4, \AM~4v/24}%
\Ballnote{This st.~is paraphrased in \Gylfaginning~39:4--5: \emph{Geit sú, er Heið·rún heitir, stendr uppi á Val-hǫll ok bítr barr af limum trés þess, er mjǫk er nafn-frę́gt, er Lę́-raðr heitir, en ór spenum hennar rennr mjǫðr sá, er hón fyllir skap-ker hvern dag.  Þat er svá mikit, at allir Ein-herjar verða full-drukknir af.} ‘That goat which is called Heathrune stands upon Walhall and bites the leaves from the branches of that tree which is very renowned, which is called Leedred.  But out of her teats flows such mead that she fills the shape-vat every day.  It is so great that all the Oneharriers become fully drunk thereof.’}%
\alst{H}ęið·ru̇n hęitir gęit, \hld\ es stęndr \edtrans{\alst{h}ǫllu ȧ (Hęrja-fǫðrs)}{on the hall of the Father of Hosts}{\Bfootnote{On Walhall.  \emph{Hęrja-fǫðrs} ‘of the Father of Hosts’ is almost certainly an unmetrical later insert in order to clarify which hall the poet is referring to and has therefore been placed in parentheses.}} &
\ind ok bítr af \edtrans{\alst{L}ę́·raðs}{Leered}{\Bfootnote{Ugdrassle’s Ash, as is seen by the mention of the stag in st.~26/1 and again in st.~36/3.}} \alst{l}imum; &
\edtrans{\alst{sk}ap-kęr}{shape-vat}{\Bfootnote{According to \CV\ the central beer-vat, from which drinks were poured into smaller vessels.}} fylla \hld\ skal \edtrans{hins \alst{sk}íra mjaðar}{the pure mead}{\Bfootnote{The goat’s milk is mead.}}, &
\ind kná⸗at sú \alst{v}ęig \alst{v}ana⸗sk.\eva

\bvb Heathrune is the goat called which stands on the hall (of the Father of Hosts) \\
\ind and bites from Leered’s branches. \\
The shape-vat she shall fill with the pure mead; \\
\ind those draughts do not wane.\evb\evg


\bvg\bva\mssnote{\Regius~10r/6, \AM~4v/26}%
\Ballnote{\Gylfaginning~39:7--8 continues with its paraphrase: \emph{Þá mę́lti Hárr: „Enn er meira mark at of hjǫrt’inn Eik·þyrni, er stendr á Val-hǫll ok bítr af limum þess trés, en af hornum hans verðr svá mikill dropi, at niðr kemr í Hver-gelmi,} ‘Then spoke High: “Yet a greater feature of it [viz.~Walhall] is the hart Oakthirner, who stands upon Walhall and bites from the branches of that tree.  But from his horns comes such great dripping that it goes down into Wharyelmer,’}%
Ęik·þyrnir hęitir \alst{h}jǫrtr \hld\ es stęndr \edtext{\edtext{\emph{\alst{h}ǫllu}}{\Afootnote{\emph{‘hællu’} \AM}} \emph{ȧ}}{\Afootnote{emend.; rev. \Regius\AM}} (Hęrja-fǫðrs) &
\ind ok bítr af \alst{L}ę́·raðs \alst{l}imum; &
ęn af hans \alst{h}ornum \hld\ drýpr ï \edtext{\alst{H}ver-gęlmi}{\Afootnote{\emph{†hvergælgji†} (for *\emph{‘hvergælmi’}) \AM}} &
\ind þaðan ęiga \alst{v}ǫtn ǫll \alst{v}ega:\eva

\bvb Oakthirner is the hart called which stands on the hall (of the Father of Hosts) \\
\ind and bites from Leered’s branches, \\
but from his horns it drip into Wharyelmer; \\
\ind thence all waters have their courses:\evb\evg


\bvg\bva\mssnote{\Regius~10r/9, \AM~4v/28}%
\Ballnote{Sts.~27--28 list numerous rivers, many of which are undoubtedly mythological.  For the sake of having somewhat more natural Anglicized renditions I have taken many of the names from \textcite[780]{CleasbyVigfusson}, who identified them with real rivers in Scotland and England, although it must be said that many or all of their identifications are entirely unverifiable and probably incorrect. ---

\Gylfaginning~39:9--10 immediately continues with its paraphrase (cf.~sts.~25--26 above), clearly deriving from sts.~27--28: \emph{ok þaðan af falla þę́r ár, er svá heita: Síð, Víð, Sekin, Ekin, Svǫl, Gunn-þró, Fjǫrm, Fimbul-þul, Gípul, Gǫpul, Gǫmul, Geir-vimul. Þessar falla um ása byggðir.  Þessar eru enn nefndar: Þyn, Vín, Þǫll, Hǫll, Grǫ́ð, Gunn-þráin, Nyt, Nǫt, Nǫnn, Hrǫnn, Vína, Veg-svinn, Þjóð-numa.} ‘and thereof fall the rivers which are called so: Side, Wide, Shin, Ecken, Swale, Guththrew, Ferm, Fimblethule, Gipple, Gapple, Gamble, Garwimble.  These fall around the settlements of the Eese.  These are yet named: Tyne, Ween, Thuil, Huil, Gread, Guththrown, Nit, Nat, Nan, Ran, Wine, Wayswith, Thednumb.’
A similar list of rivers likewise deriving from sts.~27--28 is found in \Gylfaginning~4:3--4: \emph{Þá mę́lti Jafn·hár: „Fyrr var þat mǫrgum ǫldum en jǫrð var skǫpuð er Nifl·heimr var gerr, ok í honum miðjum liggr bruðr sá, er Hver·gelmir heitir, ok þaðan af falla þę́r ár, er svá heita: Svǫl, Gunn-þrá, Fjǫrm, Fimbul-þul, Slíðr ok Hríð, Sylgr ok Ylgr, Víð, Leiptr.  Gjǫll er nę́st Hel-grindum.“} ‘Then spoke Evenhigh: “It was many ages before the earth was created that Nivelham was made, and in the middle of it lies the well which is called Wharyelmer, and thereof fall the rivers which are called so: Swale, Guththrew, Ferm, Fimblethule, Slide and Ride, Sellow and Wellow, Wide, Lafter. Yell is closest to the gates of Hell.”’  To distinguish the variants the manuscript is followed by a subscript number: 1 for the paraphrase in \Gylfaginning~4:3--4, 2 for 39:9--10, so that e.g.~\Upsaliensis\textsubscript{2} refers to the form of a name in \Gylfaginning~39:9--10 in \Upsaliensis.  TODO: Incorporate these variants into the edition.}%
\alst{S}íð ok Víð, \alst{S}ę́kin ok Ęikin, \hld\ \alst{S}vǫl ok Gunn·þró, &
\ind \alst{F}jǫrm ok \alst{F}imbul·þul, &
\ind \alst{R}ïn ok \alst{R}innandi, &
\alst{G}ipul ok \alst{G}ǫpul, \hld\ \alst{G}ǫmul ok \alst{G}ęir·vimul, &
\ind þę́r \alst{h}verfa umb \alst{h}odd goða, &
\alst{Þ}yn ok Vin, \hld\ \alst{Þ}ǫll ok Hǫll, &
\ind \alst{G}rǫ́ð ok \alst{G}unn-þorin.\eva

\bvb Side and Wide, Shin and Oaken, Swale and Guththrew, \\
\ind Ferm and Fimblethule, \\
\ind Rhine and Rinning, \\
Gipple and Gapple, Gamble and Garwimble--- \\
\ind they circle round the hoard of the Gods \ken*{= Osyard}--- \\
Tyne and Ween, Thuil and Huil, \\
\ind Gread and Guththorn.\evb\evg


\bvg\bva\mssnote{\Regius~10r/12, \AM~5r/1}%
\edtext{\alst{V}ina}{\Afootnote{\emph{‘Víná’} \Regius \emph{;} \emph{†Vma†} (for *\emph{Vina}) \AM}} hęitir ęnn, \hld\ ǫnnur \alst{V}eg-svinn, &
\ind \alst{þ}riðja \alst{Þ}jóð-numa; &
\alst{N}yt ok \alst{N}ǫt, \hld\ \alst{N}ǫnn ok Hrǫnn, &
\alst{S}líð ok Hríð, \hld\ \alst{S}ylgr ok Ylgr, &
\alst{V}íð ok \edtext{\alst{V}ǫ̇n}{\Afootnote{\emph{Vǫ́ð} \AM}}, \hld\ \alst{V}ǫnd ok Strǫnd, &
\alst{G}jǫll ok Lęiptr; \hld\ þę́r falla \alst{g}umnum nę́r &
\ind \edtext{en}{\Afootnote{so \AM \emph{;} \emph{es} \Regius}} falla til \alst{h}ęljar \alst{h}eðan. \eva

\bvb Wine is yet another called, another Wayswith, \\
\ind a third Thednumb; \\
Nit and Nat, Nan and Ran, \\
Slide and Ride, Sellow and Wellow, \\
Wide and Ween, Wand and Strand, \\
Yell and Lafter---they fall near to men, \\
\ind but fall hence to Hell.\evb\evg


\bvg\bva\mssnote{\Regius~10r/15, \AM~5r/4, \\
\RegiusProse~5r/1, \Trajectinus~TODO, \\
\Wormianus~TODO, \Upsaliensis~TODO}
\Ballnote{Thunder’s movement is so powerful that he cannot ride across the rainbow bridge as this would cause it to burn up; for the violent movement of this god cf.~\textlink{Thrymskvida}[21], \textlink{Lokasenna}[55]. --- St.~29 is cited in \Gylfaginning~15:21, the chapter about how the Gods get to their Thing, introduced by the words: \emph{En Þórr gengr til dóms’ins ok veðr ár þę́r, er svá heita:} ‘But Thunder walks to the Judgment and wades through the rivers which are called so:’.}%
\edtrans{\edtext{\alst{K}ǫrmt ok Ǫrmt}{\Afootnote{\emph{†kaurnit ok aurnit†} \Trajectinus}}}{Carmet and Armet}{\Bfootnote{Two interesting river-names.  \emph{Kǫrmt} ‘Carmet’ is the well-attested ON name of the island Karmøy in Rogaland, Norway (e.g. ÞSjár Frag 1, ESk Lv 9, Þul \emph{Eyja} 3 in \Skp\ 3), on whose northern side the ancient royal seat Awaldsness (ON \emph{Ǫg·valds-nęs}) is located.  The present stanza, however, cannot be referring to the island itself but instead seems to refer to Karmsund, the sound separating the island from the mainland.  (For reference, at its narrowest point---the place of the Karmsund bridge---the sound is only ca.~190 m wide, while at Ogwaldsness two km south it is ca.~300 m.)

The rhyming \emph{Ǫrmt} is a hapax only otherwise found in the duplicate line in Þul \emph{Á} 6.  No descendant is attested in Norway, but \textcite[780]{CleasbyVigfusson} connect it to the river Armet in Scotland.  Still, such a geographic spread seems quite unlikely, and it is more probable that the Scottish river-name, if related at all to \emph{Ǫrmt}, is taken from that of a now-lost Norwegian sound or island in the vicinity of Carmet.

\textcite{Olsen1925} derives the names from \emph{karmr} and \emph{armr}, terms for the parallel parts of a pen in a sheep-fold, which would give them a sense like ‘low protective wall’.  \citeauthor{Olsen1925} considers the present use of the names mythological and of likely Icelandic origin, deriving from a confusion of vaguely remembered names of real locations; he argues that the use of the island-name for the sound would be offensive to a Norwegian with good geographic awareness.  And yet the likeness of the islands or their sounds to a sheep-fold fits very well geographically with the area around Karmøy, so that this argument is not entirely convincing.}}%
\hld\ ok \edtext{\alst{k}ęr-laugar tvę́r}{\lemma{Kǫrmt \dots\ tvę́r ‘Carmet \dots\ Carlays’}\Bfootnote{This line also appears in a \inx[C]{thule} or list of poetic synonyms (Þul \emph{Á} 6 in \Skp\ 3), where however \emph{ok} ‘and’ in the a-verse is replaced with the name \emph{Lęiptr} ‘Lafter’ as found in st.~28/6 above.}} &
\ind \edtrans{\alst{þ}ę́r skal \alst{Þ}ȯrr vaða}{these shall Thunder wade}{\Bfootnote{Thunder is commonly associated with wading; cf.~\textlink{Harbardsljod}[13]/3 n., \textlink{EddicFragments}[F6][VI].}} &
\edtext{\alst{d}ag hvęrn}{\Afootnote{so \AM\RegiusProse\Trajectinus\Wormianus \emph{;} rev.~\Upsaliensis \emph{;} \emph{hvęrjan dag} \Regius}} \hld\ es \alst{d}ǿma fęrr &
\ind at \alst{a}ski \alst{Y}gg·drasils; &
því’t \edtrans{\alst{Ǫ̇}s-brú}{Os-bridge}{\Bfootnote{The Rainbow Bridge, Bivrest.  Cf.~st.~45/4 below.}} \hld\ brinn \alst{ǫ}ll \edtext{loga}{\Afootnote{add.~\emph{ęn} \Upsaliensis}} &
\ind \alst{h}ęilǫg vǫtn \edtrans{\alst{h}lóa}{boil}{\Afootnote{\emph{flóa} ‘turn hot’ \Upsaliensis}\Bfootnote{A hapax, \Upsaliensis’s reading is unmetrical.  \textcite[s.v.~\textit{hlóa}]{OrdsifjaBok} connects it with the adj.~\emph{hlę́r} ‘hot, mild’, and this seems likely enough.}}.\eva

\bvb Carmet and Armet and the two Carlays, \\
\ind these shall Thunder wade \\
every day, when to judge he goes \\
\ind at \inx[P]{Ugdrassle’s Ash}, \\
for the Os-bridge \ken{rainbow} all will burn with flame; \\
\ind the holy waters boil.\evb\evg


\bvg\bva\mssnote{\Regius~10r/17, \AM~5r/6}%
\Ballnote{This stanza clearly serves as a source for \Gylfaginning~15:18--20: \emph{Hestar ása’nna heita svá: Sleipnir er batstr, hann á Óðinn; hann hefir átta fǿtr.  Annarr er Glaðr, þriði Gyllir, fjórði Glenr, fimmti Skeið·brimir, sétti Silfrin·toppr, sjaundi Sinir, átti Gísl, níundi Fal·hófnir, tíundi Gull·toppr, ellipti Létt·feti.  Baldrs hestr var brenndr með hónum.} ‘The horses of the Eese are called so: Slapner is the best; Weden owns him; he has eight legs.  The second is Glad, third Gilder, fourth Glen, fifth Sheathbrimmer, sixth Silvrentop, seventh Sinewer, eighth Yissel, ninth Fallowhooven, tenth Goldtop, eleventh Lightfeet.  Balder’s horse was burned with him.’  TODO: Incorporate variants from \GylfMS.}%
\alst{G}laðr ok \alst{G}yllir, \hld\ \alst{G}lęr ok Skęið·brimir, &
\ind \alst{S}ilfrin·toppr ok \alst{S}inir, &
\alst{G}ísl ok \edtext{Fal·hófnir}{\Afootnote{\emph{‘falæpnir’} \AM}}, \hld\ \alst{G}oll·toppr ok Lėtt·feti, &
\ind þęim ríða \alst{ę̇}sir \alst{jó}um &
\alst{d}ag hvęrn \hld\ es \alst{d}ø̇ma \edtext{fara}{\Afootnote{om. \AM}} &
\ind at \alst{a}ski \alst{Y}gg·drasils.\eva

\bvb Glad and Gilder, Glare and Sheathbrimmer, \\
\ind Silvrentop and Sinewer; \\
Yissel and Fallowhooven, Goldtop and Lightfeet; \\
\ind on these horses ride the Eese \\
every day when to judge they go \\
\ind at \inx[P]{Ugdrassle’s Ash}.\evb\evg


\bvg\bva\mssnote{\Regius~10r/20, \AM~5r/8}%
\Ballnote{Ugdrassle’s Ash, the World Tree, stands at the center of the World, being the \emph{axis mundi} around which the Cosmos revolves.  This is how we should understand its being the origin of the three roots going into the three realms.  Through these roots it is also possible to travel from the land of men to the land of the dead (and vice-versa), and this is likely significant to the myth of Weden’s self-hanging upon the Tree (\textlink{Havamal}[138]).  Note that the Gods are not mentioned, for the Ash stands beside Walhall as shown by sts.~25--26 above.}%
\alst{Þ}ríar rǿtr \hld\ standa ȧ \alst{þ}ría vega &
\ind \edtext{undan}{\Afootnote{\emph{und} \AM}} \alst{a}ski \alst{Y}gg·drasils; &
\alst{H}ęl býr \edtrans{und}{enclosed by}{\Bfootnote{ON \emph{undir} (short form \emph{und}) ‘under’ can also refer to something enclosed, e.g. by a high wall, so that the residents of a walled city are said to live “under” that city.  Since the roots are prob.~not thought to be literally overhead above the listed worlds, this is the way the preposition ought to be translated.}} ęinni, \hld\ annarri \alst{h}rïm-þursar, &
\ind þriðju \alst{m}ęnnskir \alst{m}ęnn.\eva

\bvb Three roots grow in three directions \\
\ind from beneath Ugdrassle’s Ash. \\
Hell lives enclosed by one, by the other the \inx[P]{Rime-Thurses}, \\
\ind by the third manly men.\evb\evg


\bvg\bva\mssnote{\Regius~10r/22, \AM~5r/9}%
\Ballnote{Sts.~32--36 describe the various animals around Ugdrassle’s Ash.  Religiously we ought to understand that the World Tree is under attack by the forces of evil and entropy, chief among them Nithehewer, the corpse-eating dragon of \textlink{Voluspa}[38], 62. --- This st.~is paraphrased in \Gylfaginning~16:1--5: \emph{Þá mę́lti Gangleri: „Hvat er fleira at segja stór-merkja frá aski’num?“ Hár segir: „Mart er þar af at segja. Ǫrn einn sitr í limum asks’ins, ok er hann margs vitandi, en í milli augna honum sitr haukr sá, er heitir Veðr-fǫlnir. Íkorni sá, er heitir Rata-toskr, rennr upp ok niðr eptir ask’inum ok berr ǫfundar orð millum arnar’ins ok Níð·hǫggs.} ‘Gangler spoke: “What more great signs are there to be told from the ash?” High says: “There is much to tell about it.  An eagle sits in the branches of the ash and he is much knowing; but between his eyes sits the hawk called Weatherfalner.  The squirrel which is called Wratetush runs up and down along the ash and carries words of insult between the eagle and Nithehewer.”’}%
\alst{R}ata-toskr hęitir íkorni \hld\ es \alst{r}inna skal &
\ind at \alst{a}ski \alst{Y}gg·drasils; &
\alst{a}rnar \alst{o}rð \hld\ hann skal \alst{o}fan bera &
\ind ok sęgja \alst{N}íð·hǫggvi \alst{n}iðr.\eva

\bvb Wratetusk is the squirrel called which shall run \\
\ind along Ugdrassle’s Ash. \\
The eagle’s words shall he carry from above, \\
\ind and tell Nithehewer below.\evb\evg


\bvg\bva\mssnote{\Regius~10r/23, \AM~5r/11}%
\Ballnote{This st.~is paraphrased in \Gylfaginning~16:6, immediately following the paraphrase of the last st.: \emph{En fjórir hirtir renna í limum asks’ins ok bíta barr; þeir heita svá: Dáinn, Dvalinn, Dún-eyrr, Dura-þrór.} ‘But four harts run in the limbs of the ash and bite its leaves; they are called thus: Dowen, Dwollen, Downeer, Doorthrew.’}%
\alst{H}irtir ’ro ok fjórir \hld\ þęir’s af \edtext{\alst{h}ę́fingar}{\Afootnote{\emph{hę́fingjar} \AM}} ȧ &
\ind \edtrans{\edtext{\alst{g}ag-halsir}{\Afootnote{so \AM \emph{;} \emph{‘á gag-halsir’} \Regius}} \edtext{\alst{g}naga}{\Afootnote{\emph{ganga} \AM}}}{with turned necks gnaw}{\Bfootnote{An iconographic parallel to this line presents itself in the common Wiking Age depictions of great horned beasts turning their necks backwards, e.g. on the runestone E 2 or the round brooch from Espinge, Scania (HST 108124).}}: &
\alst{D}áinn ok \alst{D}valinn, \hld\ \edtext{\alst{D}u̇n·ęyrr}{\Afootnote{\emph{Dyn·ęyrr} \AM}} ok \edtext{\alst{D}ura·þrór}{\Afootnote{\emph{Dyra·þrór} \AM}}.\eva

\bvb There are four harts, too, who from the buds thereon \\
\ind with turned necks gnaw: \\
Dowen and Dwollen, Downeer and Doorthrew.\evb\evg


\bvg\bva\mssnote{\Regius~10r/25, \AM~5r/12, \\
\RegiusProse~5r/24, \Trajectinus~TODO, \\
\Wormianus~TODO, \Upsaliensis~TODO}%
\Ballnote{Sts.~34--36 are cited in \Gylfaginning~16 in the order 36, 34, 35.  34--35 are treated as a single stanza.  St.~36 is introduced by \Gylfaginning~16:7--8: \emph{En svá margir ormar eru í Hvergelmi með Níð·hǫgg, at engi tunga má telja; svá segir hér:} ‘But there are so many wyrms in Wharyelmer with Nithehewer that no tongue may tally them.  So it says here:’  Between 36 and 34--35 \Gylfaginning~16:9 reads: \emph{Svá er enn sagt:} ‘Further it is said so:’.}%
\alst{O}rmar flęiri \hld\ liggja und \alst{a}ski \alst{Y}gg·drasils &
\ind an \edtext{þat}{\Afootnote{om. \AM}} of \alst{h}yggi \alst{h}vęrr &
\ind \alst{ȯ}·sviðra \alst{a}pa:\eva

\bvb More wyrms lie under Ugdrassle’s Ash \\
\ind than any one would think \\
\ind among unwise \inx[C]{ape}[apes]:\evb\evg


\bvg\bva\mssnote{\Regius~10r/26, \AM~5r/13, \\
\RegiusProse~5r/25, \Trajectinus~TODO, \\
\Wormianus~TODO, \Upsaliensis~TODO}%
\edtext{\alst{G}óinn ok Móinn}{\Afootnote{\emph{Góni ok Móni} \Upsaliensis}}, \hld\ þęir ’ro \alst{G}raf·vitnis \edtrans{synir}{sons}{\Afootnote{\emph{liðar} ‘comrades’ \Upsaliensis}}, &
\ind \alst{G}rá·bakr ok \edtext{\alst{G}raf·vǫlluðr}{\Afootnote{\emph{Grap·vǫlluðr} \Wormianus}}, &
\edtext{\alst{Ó}fnir}{\Afootnote{\emph{Ópnir} \AM}} ok Sváfnir, \hld\ hygg’k at \alst{ę́} skyli &
\ind \edtext{\alst{m}ęiðs}{\Afootnote{\emph{męiðar} \AM \emph{;} \emph{†meðs†} \RegiusProse\Trajectinus}} \edtext{kvistu}{\Afootnote{\emph{kvistum} \RegiusProse\Trajectinus\Wormianus \emph{;} \emph{kostum} \Upsaliensis}} \alst{m}áa.\eva

\bvb Gowen and Mowen---they are Gravewitner’s sons--- \\
\ind Greyback and Gravewalled. \\
Ovner and Swebner I think always shall \\
\ind maim the branches of the tree.\evb\evg


\bvg\bva\mssnote{\Regius~10r/28, \AM~5r/14, \\
\RegiusProse~5r/22, \Trajectinus~TODO, \\
\Wormianus~TODO, \Upsaliensis~TODO}%
\alst{A}skr \alst{Y}gg·drasils \hld\ drýgir \alst{ę}rfiði &
\ind \alst{m}ęira an \alst{m}ęnn \edtext{of}{\Afootnote{so \AM\Upsaliensis \emph{;} om. \Regius\RegiusProse\Trajectinus\Wormianus}} viti: &
\alst{h}jǫrtr bítr \edtext{ofan}{\Afootnote{\emph{neðan} (corruption from l.~4) \Upsaliensis}} \hld\ ęn ȧ \alst{h}liðu fu̇nar, &
\ind skęrðir \alst{N}íð·hǫggr \alst{n}eðan.\eva

\bvb Ugdrassle’s Ash suffers hardship \\
\ind more than men might know: \\
A hart bites it from above and it rots on the side; \\
\ind Nithehewer harms it from below.\evb\evg


\bvg\bva\mssnote{\Regius~10r/30, \AM~5r/16, \\
\RegiusProse~9r/4, \Trajectinus~TODO, \\
\Wormianus~TODO, \Upsaliensis~TODO}%
\Ballnote{St.~37 is cited following \Gylfaginning~36:2.  \Gylfaginning~36:1--2 introduces it: \emph{Enn eru þę́r aðrar er þjóna skulu í Val-hǫll, bera drykkju ok gę́ta borð-búnaðar ok ǫl-gagna.  Svá eru þę́r nefndar í Grímnis-mǫ́lum:} ‘Yet there are other [heavenly] women who shall serve in Walhall, carry the drinks and prepare the table-cloth and ale-stuffs.  So are they named in the Speeches of Grimner:’.}%
\alst{H}rist ok Mist \hld\ vil’k at mér \alst{h}orn beri, &
\ind \edtext{\alst{Sk}eggj·ǫld}{\Afootnote{so \Regius\RegiusProse\Trajectinus\Wormianus \emph{;} \emph{Skęgg·ǫld} \AM\Upsaliensis.}} ok \alst{Sk}ǫgul, &
\edtrans{\alst{H}ildr ok Þrúðr}{Hild and Thrith}{\Afootnote{\emph{Hildi ok Þrúði} \Regius}\Bfootnote{The \Regius\ reading stems from *\emph{ꝺꝛ, *ðꝛ} with \emph{r} rotunda being mistakenly copied as \emph{ꝺı, ðı} and is certainly not an independent oral variant of these well-attested walkirrie names.  This is apparent upon viewing the facsimile images.}}, \hld\ \alst{H}lǫkk ok \edtrans{\alst{H}ęr·fjǫtur}{Harfetter}{\Afootnote{\emph{Hęr·fjǫtra} \Upsaliensis}}, &
\ind \edtext{\alst{G}ǫll}{\Afootnote{\emph{Gjǫll} \Upsaliensis}} ok \edtrans{\alst{G}ęira·hǫð}{Garhath}{\Afootnote{so \GylfMS \emph{;} \emph{Gęir·ǫlul} ‘Garalel’ \Regius \emph{;} \emph{Gęir·rǫmul} ‘Garramble’ \AM. --- add.~\emph{ok} ‘and’ \Upsaliensis.}}, &
\edtext{\alst{R}and·gríð}{\Afootnote{\emph{Rand·gríðr} \Upsaliensis}} ok \edtext{\alst{R}áð·gríð}{\Afootnote{\emph{Ráð·gríðr} \Upsaliensis}}, \hld\ \edtext{ok}{\Afootnote{om. \Regius}} \alst{R}ęgin·lęif; &
\ind \edtrans{þę́r bera \alst{ęi}n-hęrjum \edtext{\alst{ǫ}l}{\Afootnote{\emph{ǫl} corrected from \emph{ǫll} in same hand \AM}}}{they bring the Oneharriers ale}{\Bfootnote{As cupbearers in Walhall.  In Iron Age Germanic society pouring drinks at a feast was traditionally done by the ruler’s kinswomen, the most famous literary example being Rothgar’s wife and daughter in \Beowulf; see \textcite{Enright1996} and \textcite{Riseley2014}.
The Walkirries may have been thought of as the daughters of Weden; see note to \textlink{Voluspa}[30]/5.  For the reception of dead warriors see also note to st.~53/3 below.}}.\eva

\bvb Rist and Mist I would have bring me a horn--- \\
\ind Shedgeld and Shagle, \\
Hild and Thrith, Lank and Harfetter, \\
\ind Gall and Garhath, \\
Randgrith and Redegrith and Rainlaf--- \\
\ind they bring the Oneharriers ale.\evb\evg


\bvg\bva\mssnote{\Regius~10r/32, \AM~5r/18}%
\Ballnote{Sts.~38--40 deal with phenomena concerning the sun.
The description of the sun as being pulled by a horse-driven chariot is attested in literature and archeology going back to the Bronze Age.  A notable archeological find from the Germanic cultural area is the ca.~1400 BC Trundholm sun chariot found on Zealand, Denmark.  This cultic object belonging to the Nordic Bronze Age features a large disc placed on a six-wheeled chariot pulled by a horse (figg.~\ref{fig:trundholm-front}, \ref{fig:trundholm-back}, page~\pageref{fig:trundholm-front}).  Only one side of the disc is gilt and so the chariot may also depict the horses pulling the Day and Night, for which see \textlink{Vafthrudnismal}[11]--14.
Indo-European literary parallels for horses pulling the sun include the Vedic \Rigveda\ 1.50.8--9, 4.13.3, 5.45.9 and the Greek Homeric Hymn to \textgreek{Ἥλῐος}.

This st.~is referenced in \Gylfaginning~11:4--7, about the origin of the Sun and Moon: \emph{En goð’in reiddu⸗sk þessu of·drambi tóku þau systkin ok settu upp á himin, létu Sól keyra þá hesta, er drógu kerru sólar’innar, þeirar er goð’in hǫfðu skapat til at lýsa heimana af þeiri síu, er flaug ór Múspells-heimi.  Þeir hestar heita svá, Ár·vakr ok Al·sviðr. En undir bógum hesta’nna settu goð’in tvá vind-belgi at kǿla þá, en í sumum frǿðum er þat kallat ísarn-kol.} ‘But the Gods were angered by this arrogance and took the two siblings and placed them up in the heaven; they let Sun drive the horses that pulled the chariot of the sun which the Gods had created to brighten the Homes from the sparks which flew out of Muspellsham.  Those horses are so named: Yorewaker and Allswith.  But under the shoulders of the horses the Gods placed two air bellows to cool them, but in some sources [presumably this st.] those are called iron-cooling.’
Cf.~\textlink{Vafthrudnismal}[23] n., which cites \Gylfaginning~11:1--3.}%
\edtrans{\alst{Á}r·vakr ok \alst{A}l·sviðr}{Yorewaker and Allswith}{\Bfootnote{The two horses which pull the sun-chariot also appear in \textlink{Sigrdrifumal}[15]/2; cf.~note to the next st.}}, \hld\ þęir skulu \alst{u}pp heðan &
\ind \edtrans{\alst{s}vangir}{famished}{\Bfootnote{M.~nom.~pl.~of \emph{svangr} ‘hungry, starved’.  This adj.~is used particularly of horses exhausted after long exertion, for which cf.~\textlink{Oddrunargratr}[3]/4, \textlink{HelgakvidaOne}[42]/4, Sigv \emph{Austv} 11 (in \Skp\ 1).}} \alst{s}ól draga; &
ęn und þęira \alst{b}ógum \hld\ fǫ́lu \alst{b}líð ręgin, &
\ind \alst{ę̇}sir, \alst{í}sarn-kol.\eva

\bvb Yorewaker and Allswith---they shall upward hence, \\
\ind famished, pull the sun, \\
but under their shoulders the blithe Reins hid, \\
\ind the Eese, iron-cooling.\evb\evg


\bvg\bva\mssnote{\Regius~10v/2, \AM~5r/20}%
\Ballnote{The sun-disc was apparently thought to be a translucent shield which protected the earth from the full power of the sun behind it.  Without it the “crags and surf” (\textsc{land} and \textsc{sea}; the totality of the earth) would burn.  Cf.~\textlink{Sigrdrifumal}[15]/1 which mentions the “shield that stands before the shining god \ken{sun}” and is probably drawing on the present stanza.}%
\edtext{\alst{S}valinn}{\Afootnote{so (\emph{‘svalin’}) \AM \emph{;} \emph{Svǫl} \Regius}\Bfootnote{The \AM\ reading is preferred since the shield is described by the m.~\emph{hann} in both mss., whereas \Regius’s \emph{Svǫl} looks like a feminine.  \emph{Svalinn} is confirmed by its inclusion in the poetic list of synonyms for shields, Þul \emph{Skjaldar} 2/3 (in \Skp\ 3).}} hęitir, \hld\ hann stęndr \alst{s}ólu fyrir, &
\ind \alst{sk}jǫldr \alst{sk}ïnanda goði; &
\alst{b}jǫrg ok \alst{b}rim \hld\ vęit’k at \alst{b}rinna skulu, &
\ind ef hann \alst{f}ęllr ï \alst{f}rȧ.\eva

\bvb Swalen one is called; it stands before the sun--- \\
\ind a shield before the shining god \ken{sun}. \\
Crags and surf I know shall burn, \\
\ind if it falls away.\evb\evg


\bvg\bva\mssnote{\Regius~10v/4, \AM~5r/21}%
\Ballnote{For discussion on these wolves see \textlink{Voluspa}[40] n. --- This stanza is closely paraphrased in \Gylfaginning~12:4--5, following a question about why the sun hurries so quickly: \emph{Hárr segir: „Þat eru tveir ulfar, ok heitir sá, er eptir henni ferr, Skoll.  Hann hrę́ði⸗sk hón, ok hann mun taka hana.  En sá heitir Hati Hróð·vitnis-son, er fyrir henni hleypr, ok vill hann taka tungl’it, ok svá mun verða.“} ‘High says: “It is due to the two wolves, and the one who runs behind her is called Sholl.  She fears him, and he would catch her.  But the one who runs ahead of her is called Hater Rothwitner’s son, and he wants to catch the moon, and so it will remain.”’}%
\alst{Sk}oll hęitir ulfr, \hld\ es fylgir hinu \alst{sk}ír-lęita &
\ind goði \edtrans{til \edtext{\alst{v}arna}{\Afootnote{\emph{†vaurna†} \AM}} \alst{v}iðar}{to the shelter of the woods}{\Bfootnote{The sun is said to set in the forest.  This is an interesting geographic detail since it points away from most of coastal (western) Norway, where the sun sets in the sea, and Iceland, which does not have any expansive forests.  The description would however fit well for the Oslofjord area.}}, &
ęn annarr \edtext{\alst{H}ati}{\Afootnote{add.~\emph{hann ’s} ‘he is’ \Regius.}} \hld\ \alst{H}róð·vitnis sonr, &
\ind sá skal fyr \alst{h}ęiða brúði \alst{h}imins.\eva

\bvb \inx[P]{Sholl} is the wolf called which follows the pure-faced \\
\ind god \ken{sun} to the shelter of the woods, \\
but the other one is \inx[P]{Hater}, \inx[P]{Rothwitner}’s son; \\
\ind he shall [run] ahead of the clear bride of heaven \ken{sun}.\evb\evg


\bvg\bva\mssnote{\Regius~10v/6, \AM~5r/23, \\
\RegiusProse~3r/30, \Trajectinus~4r/10, \\
\Wormianus~TODO, \Upsaliensis~TODO, \\
\AMb~9v/14, \EddaBms~3v/11}%
\Ballnote{St.~41 is clearly closely related to \textlink{Vafthrudnismal}[21]; for the Creation see there and \textlink{Voluspa}[3]--4. ---
Sts.~41--42 are cited in sequence following \Gylfaginning~8:20--25, which describes the Earth: \emph{[...] Hón er kring-lótt útan, ok þar útan um liggr inn djúpi sjár, ok með þeiri sjávar-strǫndu gǫ́fu þeir lǫnd til byggðar jǫtna ę́ttum.  En fyrir innan á jǫrðu’nni gerðu þeir borg um-hverfis heim fyrir ó·friði jǫtna.  En til þeirar borgar hǫfðu þeir brár Ymis jǫtuns ok kǫlluðu þá borg Mið·garð.  Þeir tóku ok heila hans ok kǫstuðu í lopt ok gerðu af ský’in.  Svá sem hér segir.} ‘It is round without and surrounding it there lies the deep sea; and along that sea-shore they [= the Gods] gave the land to the settlement by the lineages of the Ettins.  But further within upon the Earth they made a stronghold around the Home against the enmity of the ettins, but for that stronghold they used the eyebrows of the ettin Yimer and they called the stronghold Middenyard.  They also took his brains and cast them into the air and thereof made the clouds, just as it says here:’ --- Sts.~41--42 are also cited with numerous variants in \emph{Lítla Skálda}, a short text on Scaldic poetics that likely served as a source for Snorre’s Edda (cf.~\textlink{EddicFragments}[F8][VIII]).}%
Ór \alst{Y}mis holdi \hld\ vas \alst{jǫ}rð of skǫpuð, &
\ind ęn ór \edtrans{\alst{s}vęita}{his blood}{\Afootnote{\emph{hans sára svęita} ‘the blood of his wounds’ \AMb\EddaBms}} \edtrans{\alst{s}jór}{sea}{\Afootnote{so \AM\Upsaliensis\AMb\EddaBms \emph{;} \emph{sę́r} ‘id.’ \Regius \emph{;} \emph{sjár} ‘id.’ \RegiusProse\Trajectinus\Wormianus.}}, &
\edtext{\edtext{\alst{b}jǫrg ór \edtext{\alst{b}ęinum}{\Afootnote{\emph{bęinu’num} \Trajectinus}}, \hld\ \edtext{\alst{b}aðmr}{\Afootnote{\emph{†baðrmr†} \AM}}}{\lemma{bjǫrg \dots\ baðmr ‘mountains \dots\ woods’}\Bfootnote{Formulaic alliterative pair describing the natural world in a context of creation or destruction, also found in the OHG \textlink{Wessobrunn} 3 (\emph{paum noh perek}) and \textlink{Muspilli} 50 (\emph{perga, poum}).}} ór hári, &
\ind ęn \edtrans{ór \alst{h}ausi \alst{h}iminn}{from his skull the heaven}{\Afootnote{\emph{himinn ór hausi hans} ‘the heaven from his skull’ \AMb\EddaBms}}.}{\lemma{bjǫrg \dots\ himinn.}\Afootnote{Abbrev.~as \emph{b. or. b. b. or. h. En or. h. h.} \Upsaliensis \emph{,} cf.~\textcite[137, 145--146]{MartenssonPalsson2008}.}}\eva

\bvb From \inx[P]{Yimer}’s flesh was the Earth created, \\
\ind but from his blood the sea, \\
mountains from his bones, woods from his hair, \\
\ind but from his skull the Heaven.\evb\evg


\bvg\bva\mssnote{\Regius~10v/8, \AM~5r/25, \\
\RegiusProse~3r/31, \Trajectinus~4r/11, \\
\Wormianus~TODO, \Upsaliensis~TODO, \\
\AMb~9v/16, \EddaBms~3v/12}%
\edtext{Ęn ór hans \alst{b}rǫ́um \hld\ gørðu \alst{b}líð ręgin &
\ind \alst{M}ið·garð \alst{m}anna sonum,}{\lemma{Ęn ór hans brǫ́um \hld\ gørðu blíð ręgin / Mið·garð manna sonum ‘But from his brows the blithe Reins made Middenyard for the sons of men’}\Bfootnote{The Gods fenced in Middenyard (‘the Middle Enclosure’) by using the strands of Yimer’s eyebrows as poles, as explained in \Gylfaginning~8:22--23.  The creation of this enclosure was a gift from the kind Gods to mankind in order to protect us from the forces of destruction without; cf.~\textlink{Voluspa}[16]/2 n.}} &
\edtext{ęn}{\Afootnote{\emph{ok} \Upsaliensis}} ór hans \alst{h}ęila \hld\ vǫ́ru þau \edtext{hin}{\Afootnote{om. \Trajectinus}} \edtrans{\alst{h}arð-móðgu}{hard-minded}{\Afootnote{\emph{hríð-fęldu} ‘stormy’ \AMb\EddaBms}} &
\ind \alst{sk}ý ǫll of \alst{sk}ǫpuð.\eva

\bvb But from his brows the blithe \inx[P]{Reins} made \\
\ind \inx[P]{Middenyard} for the sons of men, \\
but from his brains were the hard-minded \\
\ind clouds all created.\evb\evg


\bvg\bva\mssnote{\Regius~10v/9, \AM~5r/26}%
\Ballnote{This st.~is one of the most difficult in the poem and many interpretations have been made.  The traditional view (e.g. \textcite{FinnurEdda}, Bellows, Sijmons and Gering p. 208) relates it to the poem’s frame narrative.  Weden, bound between the two fires, cryptically asks for a cauldron hanging above him from the roof to be moved aside so that the Gods will be able to see him through the smoke-vent and rescue him.  This explanation, however, leaves much unexplained, namely the stanza’s placement in the gnomic wisdom section of the poem, the invocation of the obscure god Wolder, the lack of mention of any cauldron elsewhere in the poem, and the question of why the gods would bestow their grace unto the person who first set the fire which is presently torturing Weden (assuming \emph{tękr ȧ funa} means ‘kindles the fire’; see note to that phrase).

I find the interpretation of \textcite{Nordberg2005} more convincing.  Nordberg argues that the st.~is another piece of gnomic wisdom, referring to the cooking of the sacrificial meal in large cauldrons during the \inx[C]{bloot}[Bloot].  The importance of cauldrons is seen in the common name second element \emph{-kętill} or \emph{-kęll} (cognate with English “kettle”) and has textual support, e.g. in \HakonarSaga\ 14, describing the traditional bloot in the Throndlaw (\emph{Þrǿnda-lǫg}), Norway: \emph{At veitslu þeiri skyldu allir menn ǫl eiga; þar var ok drepinn alls konar smali ok svá hross, [...] en slátr skyldi sjóða til mann-fagnaðar; eldar skyldu vera á miðju gólfi í hofi’nu ok þar katlar yfir.} ‘At that gathering all men were to have ale; thereat also was every kind of small cattle slain and likewise horses, [...] and the fresh meat should be cooked for the enjoyment of men.  There should be fires in the middle of the floor in the hove and kettles above them.’
According to this view, the stanza is speaking of the Heavenly favour (\emph{hylli}) earned by the ritualist who kindles the cooking fire, since that act enables the Gods to become guests at the ritual meal.

Nordberg’s interpretation becomes especially interesting when one considers the immediately preceding stanzas 41--42, which describe the ordering of the world by the Gods through the killing and dismembering of Yimer, the primordial victim.  That the killing of Yimer was in fact a sacrifice is supported by the manner in which it was done: beheading, which was the typical manner of slaying sacrificial bulls in the Wiking Age (cf.~\textlink{Vafthrudnismal}[21]/4 n., \cite{Lucas2007}).
In other Indo-European religions the first sacrifice of a Great Being (most famously the Vedic \textoi{Púruṣa}, \Rigveda\ 10.90) serves as the model for all future sacrifice, the performance of which reenacts the Creation and thereby re-establishes the continued existence of the World and upholding the social order (\cite[\emph{passim}]{Lincoln1986}; cf. \textlink{Voluspa}[7]/2 n.).  That the purpose of the sacrifice is the re-establishment of Creation is said explicitly in the ritual language of the Vedic fire-sacrifice \parencite[11, 78--80]{EliadeReturn}.
With this in mind we may consider how the sequence \textlink{Grimnismal}[41]--43 likely forms a unit: sts.~41--42 describe the Creation of the World through the First Bloot (sacrifice) in the mythic past (\emph{in illo tempore}) and st.~43 connects it with the Bloot in the present.}%
\edtrans{\alst{U}llar}{Wolder’s}{\Bfootnote{It is uncertain why the obscure god Wolder is invoked here but it cannot be simply for the sake of alliteration, since the much more well-known \emph{Óðins} ‘Weden’s’ would work just as well.  It is possible that Wolder had a particular role in the setting of the ritual fire, something which finds support in the large number of fire-striker-shaped iron amulet-rings dating to ca.~660--780 \ce\ found at the archeological site of \emph{Lilla Ullevi} (‘Wolder’s little \inx[C]{wigh}’) in Sweden \parencite{afEdholm2009}.

In Wiking Age and High Mediæval sources Wolder plays a very minor role, but his earlier significance is proven by the impressive number of cultic places named after him \parencite[116--118]{Brink2007} and is, apart from the present stanza, also hinted at by \textlink{Atlakvida}[33]/4.}} \edtrans{hylli}{holdness}{\Bfootnote{‘Favour, loyalty, grace’.  This word, derived from the adjective \emph{hollr} ‘hold; favourable, loyal, gracious’ along with the verb \emph{hylla} ‘to make hold’ is used to refer to the grace of god(s) in both Heathen and Christian texts.
See Index: \inx[C]{hold} and \inx[C]{holdness}.}} \hld\ hęfr ok \edtrans{\alst{a}llra goða}{All Gods}{\Bfootnote{Cf.~\textlink{Sigrdrifumal}[3]--4, \textlink{Lokasenna}[11], which both hail the Gods as a collective (the former as part of a genuine prayer, the latter subversively) and \Hakonarmal\ 18, where the piety of the dead king Hathkin is shown by his being greeted by \emph{rǫ́ð ǫll ok ręgin} ‘all the Redes and \inx[P]{Reins}’.

The ritual use of this phrase proves that the pre-Christian religion did in fact believe in a form of Oneness of the Gods and their identity as a united body, something also seen in the idea of the \inx[P]{Thing of the Gods} (see \textlink{Voluspa}[6] n.), where all the Gods gather to steer the fates of the world.

Similar expressions are found in other old Indo-European religions, e.g. the Vedic \textoi{víşvē dēvā́ḥ} ‘All Gods’, to Whom are dedicated numerous hymns of \Rigveda, and the Greek \textgreek{Πάν·θειον} (Pantheon), that is, a temple dedicated to All Gods.

This idea of Godly Oneness may have been the subject of dispute in the Norse world.  \textcite[53--54]{Saxo}, 1.7.2 gives us an interesting anecdote.  After Weden departs he is usurped by the obscure \emph{Mithothin} (perhaps “With-Weden”), who reforms the cult:

  \emph{Cuius secessu Mithothyn quidam prestigiis celeber, perinde ac celesti beneficio vegetatus, occasionem et ipse fingende divinitatis arripuit barbarasque mentes novis erroris tenebris circumfusas prestigiarum fama ad cerimonias suo nomini persolvendas adduxit. Hic deorum iram aut numinum violationem confusis permixtisque sacrificiis expiari negabat ideoque eis vota communiter nuncupari prohibebat, discreta superum cuique libamenta constituens. Qui cum Othino redeunte relicta prestigiarum ope latendi gratia Pheoniam accessisset, concursu incolarum occiditur.}

  ‘A certain Mithodin, a famous illusionist, was animated at his departure as if by a kindness from heaven and snatched the chance to pretend divinity himself; his reputation for magicianship clouded the barbarians’ minds with the murk of a new superstition and led them to perform holy rites to his name. He asserted that the gods’ wrath and the profanation of their divine authority could not be expiated by confused and mingled sacrifices; so he arranged that they must not be prayed to as a group, but separate offerings (\emph{libamenta}) be made to each deity. When Odin returned, the other no longer resorted to his conjuring but went off to hide in Funen, where he was rushed upon and killed by the inhabitants.’

It is not impossible that this obviously mythologised retelling may in fact reflect an actual theological conflict or attempted religious reform within paganism.}} &
\ind hvęrr’s \edtext{tękr \alst{f}yrstr ȧ \alst{f}una}{\lemma{tękr \dots\ ȧ funa ‘kindles the fire’}\Bfootnote{Viz.~the hallowed fire above which the holy cauldrons are hung; for the role of fire in Germanic and Vedic sacrifice see \textcite{Kaliff2005}. ---
An otherwise unattested phrase which is here taken as equivalent with \emph{taka ęld} ‘kindle fire’, although this is problematic since the preposition \emph{ȧ} ‘on’ modifies the sense of the verb and the phrase \emph{taka ȧ} has a variety of idiomatic senses like ‘touch, react to, get involved in, get on’, though *\emph{taka ȧ ęld} is not attested.  Another possibility is to read the ms.~\emph{afuna} as an error for \emph{af funa} ‘from the fire’, which would yield the superior sense ‘whoever first takes [the cauldrons] from the fire’ \parencite[55]{Nordberg2005}.}}, &
því’t \edtext{\alst{o}pnir hęimar \hld\ verða}{\lemma{opnir hęima verða ‘the Homes become open’}\Bfootnote{The Bloot opens up the human world to that of the Gods above.  Similar language is used in several hymns of \Rigveda\ which pray for the Heavenly Doors (\textoi{dvā́raḥ dēvī́ḥ}) to open up, e.g.~\Rigveda\ 1.13.6, 1.142.6, 2.3.5.}} umb \alst{ȧ}sa sonum, &
\ind \edtext{þȧ’s}{\Afootnote{add.~\emph{þęir} \AM}} \alst{h}ęfja af \edtrans{\alst{h}vera}{cauldrons}{\Bfootnote{Acc. pl. of \emph{hverr}, from PGmc. *\emph{hweraz}, from PIE *\emph{kʷer-} ‘pot, vessel’ \parencite[265]{Kroonen2013}.  The Sanskrit cognate \textoi{carú} is occasionally used in reference to the vat from which the ritual drink \textoi{sóma} is drunk (\Rigveda\ 10.167.4), but a particular religious significance for the PIE root cannot be reconstructed.}}.\eva

\bvb \inx[P]{Wolder}’s \inx[C]{holdness} and that of All Gods \\
\ind has whoever first kindles the fire, \\
for the \inx[C]{Home}[Homes] become open for the sons of the Eese \ken{gods}, \\
\ind when men lift off the cauldrons.\evb\evg


\bvg\bva\mssnote{\Regius~10v/11, \AM~5r/28}%
\Ballnote{The ship Shidebladner is also described in \Gylfaginning~43, which is quite likely pulling from the present stanza.}%
\alst{Í}·valda synir \hld\ gingu ï \alst{á}r-daga &
\ind \alst{Sk}íð·blaðni at \alst{sk}apa, &
\alst{sk}ipa bętst \hld\ \alst{sk}írum Fręy, &
\ind \alst{n}ýtum \alst{N}jarðar bur.\eva

\bvb Iwald’s sons went in days of yore \\
\ind Shidebladner for to shape: \\
the best of ships for the pure Free, \\
\ind for the useful Son of Nearth.\evb\evg


\bvg\bva\mssnote{\Regius~10v/13, \AM~5r/29, \\
\RegiusProse~10r/20, \Trajectinus~TODO, \\
\Wormianus~TODO, \Upsaliensis~TODO}%
\Ballnote{St.~45 concludes Weden’s mythological wisdom with a list of the greatest things in each category.  It is quite likely that this stanza concluded with \emph{Bif·rǫst brúa} ‘Bivrest of bridges’ as a \Ljodahattr\ c-verse and that the following verses are later interpolations in oral transmission.  Unlike the previous things listed they seem more trivial.  The interpolating tendency is visible in action in \AM\ which adds a final unmetrical verse at the end of the stanza. ---

The stanza is cited following \Gylfaginning~41:6--7: \emph{En satt er þat, er þú sagðir.  Mikill er Óðinn fyrir sér.  Mǫrg dǿmi finna⸗sk til þess.  Svá er hér sagt í orðum sjalfra ásanna:} ‘But true is that which thou hast said.  Weden is formidable; many examples attest to that.  So it is said here in the words of the Eese themselves.’

For \Gylfaginning~41:1--5 see \textlink{Vafthrudnismal}[41] n.}%
\alst{A}skr \alst{Y}gg·drasils, \hld\ \edtext{hann}{\Afootnote{om. \Upsaliensis}} ’s \alst{ǿ}ðstr viða &
\ind ęn \alst{Sk}íð·blaðnir \alst{sk}ipa, &
\alst{Ó}ðinn \alst{ȧ}sa \hld\ ęn \alst{jó}a Slęipnir, &
\edtrans{\alst{B}if·rǫst}{Bivrest}{\Afootnote{so \GylfMS \emph{;} \emph{Bil·rǫst} ‘Bilrest’ \Regius\AM. --- om., but add.~above the line by same hand \AM}} \alst{b}rúa \hld\ ęn \alst{B}ragi skalda, &
\alst{H}ǫ́·brók \alst{h}auka \hld\ ęn \edtrans{\alst{h}unda \edtext{Garmr}{\Afootnote{\emph{Gramr} \RegiusProse\Trajectinus}}}{Garm of hounds}{\Afootnote{add. \emph{ęn Brimir sverða} ‘and Brimmer of swords’ \AM.}}.\eva

\bvb Ugdrassle’s Ash---it is the noblest of trees, \\
\ind and Shidebladner of ships, \\
Weden of the Eese and Slopner of steeds, \\
Bivrest of bridges and Bray of scalds, \\
Highbrook of hawks and Garm of hounds.\evb\evg


\bvg\bva\mssnote{\Regius~10v/15, \AM~5v/2}%
\Ballnote{Weden announces that he has raised his gaze and made the Gods aware of his situation; they will leave their feasting at Eagre’s hall (cf.~\textlink{Hymiskvida} and \textlink{Lokasenna}) and come to his rescue.}%
\alst{S}vipum hęfi’k \edtext{nú}{\Afootnote{om.~\AM}} ypt \hld\ fyr \alst{s}ig-tíva sonum, &
\ind við þat skal \alst{v}il-bjǫrg \alst{v}aka, &
\alst{ǫ}llum \alst{ǫ̇}sum \hld\ þat skal \alst{i}nn koma &
\ind \alst{Ę́}gis bękki \alst{ȧ} &
\ind \alst{Ę́}gis drekku \alst{a}t.\eva

\bvb My gaze I’ve now lifted up before the sons of the victory-Tews; \\
\ind by that shall the willed rescue awake! \\
All the Eese shall it bring in, \\
\ind upon Eagre’s bench, \\
\ind at Eagre’s drinking.\evb\evg


\bvg\bva\mssnote{\Regius~10v/17, \AM~5v/4, \\
\RegiusProse~6r/5, \Trajectinus~TODO, \\
\Wormianus~TODO, \Upsaliensis~TODO}%
\Ballnote{Weden begins his list of names which takes up sts.~47--52; more names are given in a second list in 56--57.  Note that the name \emph{Óðinn} ‘Weden’ itself does not appear in the first list. --- The names from sts.~47--52 and 56 (but not 57) are also found in \Gylfaginning\ 20:11, introduced by the words: \emph{ok enn hefir hann nefnt⸗sk á fleiri vega, þá er hann var kominn til Geir·røðar konungs:} ‘and he has named himself in yet more ways when he was come to visit King Garfrith:’  The names are given in the order they appear in each stanza and the stanzas themselves come in the order: 47, 48, 49, 51, 52, 50, 56.  For the purpose of the present critical edition, variant names from the list in \Gylfaginning~20:11 are noted for all relevant stanzas, but only sts.~47--49 are treated as being actually cited in \Gylfaginning.  Note that all 4 mss.~of \Gylfaginning\ append \emph{Vera-týr} to the end of the list, likely taken from st.~3 above.}%
Hét⸗umk \alst{G}rïmr, \hld\ \edtext{hét⸗umk}{\Afootnote{\emph{ok} \AM\RegiusProse\Wormianus\Upsaliensis \emph{;} om.~\Trajectinus}} \edtext{\alst{G}anglęri}{\Afootnote{so \Regius\AM\Upsaliensis \emph{;}
\emph{Gangari} \RegiusProse \emph{;} \emph{Gangi} \Trajectinus \emph{;} \emph{Ganglari} \Wormianus.}}, &
\ind \alst{H}ęrjann \edtext{ok}{\Afootnote{om. \GylfMS}} \alst{H}jalm·beri, &
\alst{Þ}ękkr \edtext{ok}{\Afootnote{om. \GylfMS}} \alst{Þ}riði, \hld\ \edtext{\alst{Þ}uðr}{\Afootnote{so \AM\GylfMS \emph{;} \emph{Þunðr} (\emph{‘Þv̅ðr’}) \Regius}} \edtext{ok}{\Afootnote{om. \GylfMS}} Uðr, &
\ind \edtext{\alst{H}ęl·blindi}{\Afootnote{so \Regius\GylfMS \emph{;} \emph{Hęr·blindr} \AM}} \edtext{ok}{\Afootnote{om. \GylfMS}} \edtext{\alst{H}ǫ́r}{\Afootnote{\emph{Hárr} \Wormianus}}.\eva

\bvb I called myself Grim; I called myself Gangler, \\
\ind Harn and Helmbearer. \\
Theck and Third, Thith and Ith, \\
\ind Hellblinder and High.\evb\evg


\bvg\bva\mssnote{\Regius~10v/19, \AM~5v/5, \\
\RegiusProse~6r/6, \Trajectinus~TODO, \\
\Wormianus~TODO, \Upsaliensis~TODO}%
\alst{S}aðr \edtext{ok}{\Afootnote{om. \GylfMS}} \alst{S}vipall \hld\ \edtext{ok}{\Afootnote{om. \GylfMS}} \alst{S}ann·getall, &
\ind \alst{H}ęr·tęitr \edtext{ok}{\Afootnote{om. \GylfMS}} \alst{H}nikarr, &
\alst{B}il·ęygr, \alst{B}ál·ęygr, \hld\ \alst{B}ǫl·verkr, Fjǫlnir, &
\edtext{\alst{G}rïmr ok}{\Afootnote{om.~\GylfMS}} \alst{G}rïmnir, \hld\ \alst{G}lap·sviðr \edtext{ok}{\Afootnote{om. \GylfMS}} Fjǫl·sviðr.\eva

\bvb Sooth and Swiple and Soothgettle, \\
\ind Hartote and Nicker, \\
Bileye, Baleeye, Baleworker, Fillner, \\
Grim and Grimner, Glapswith and Fellswith.\evb\evg


\bvg\bva\mssnote{\Regius~10v/21, \AM~5v/7, \\
\RegiusProse~6r/7, \Trajectinus~TODO, \\
\Wormianus~TODO, \Upsaliensis~TODO}%
\edtext{\alst{S}íð·hǫttr}{\Afootnote{\emph{Síð·hófr} \Trajectinus}}, \alst{S}íð·skęggr, \hld\ \alst{S}ig·fǫðr, \edtext{Hnikuðr, &
\alst{A}l·fǫðr, \edtrans{\alst{V}al·fǫðr}{Walfather}{\Afootnote{om.~\AM\GylfMS}}, \hld\ \edtext{\alst{A}t·ríðr}{\Afootnote{\emph{†ettridr†} \Trajectinus}}}{\lemma{Hnikuðr \dots\ At·ríðr}\Afootnote{\emph{At·ríðr, Hnikuðr, Al·fǫðr} \Upsaliensis}} \edtext{ok}{\Afootnote{om. \GylfMS}} Farma·týr; &
\edtext{\alst{ęi}nu nafni \hld\ hét⸗umk \edtrans{\alst{a}ldri⸗gi}{never}{\Afootnote{om.~\AM}} &
\ind síðst ek \edtext{með}{\Afootnote{om.~\AM}} \alst{f}olkum \alst{f}ór.}{\lemma{ęinu \dots\ fór. ‘by \dots\ fared.’}\Afootnote{om.~\GylfMS}}\eva

\bvb Sidehat, Sideshaw, Syefather, Nicked, \\
Allfather, Walfather, Atrider, and Farm-Tew--- \\
by a single name I never called Myself \\
\ind since among man-folk I fared.\evb\evg


\bvg\bva\mssnote{\Regius~10v/23, \AM~5v/9}%
\Ballnote{Apart from the first line, which of course refers to the present situation, the underlying stories are not known.  They presumably refer to other now-lost myths involving Weden travelling in disguise. --- The Weden-names in this stanza are found in \Gylfaginning~20:11 following those of st.~52.  They appear in the order \emph{Jalkr, Kjalarr, Viðurr, Þrór} ‘Gelding, Keller, Wither, Throo’ in all mss., indicating that ll.~4 and 5 were reversed in the version of the poem used by Snorre.}%
\alst{G}rïmni mik hétu \hld\ at \alst{G}ęir·raðar, &
\ind ęn \alst{Ja}lk at \edtext{\alst{Ǫ̇}s·mundar}{\Afootnote{This name is spelled \emph{‘ǫ́smundar’} in \Regius, an archaic scribal feature suggesting that the poem was likely put to paper in the 12th century.}}; &
ęn þȧ \edtext{\alst{K}jalar}{\Afootnote{so \Regius\GylfMS \emph{;} \emph{Jalk} \AM.}} \hld\ es ek \alst{k}jalka dró, &
\ind \edtrans{\alst{Þ}rór \alst{þ}ingum at}{Throo at Things}{\Bfootnote{Perhaps a reference to a name under which Weden was invoked at the start of Things or legal assemblies.}}, &
\ind \edtrans{\alst{V}iður \emph{\alst{v}ígum at}}{Wither in battles}{\Afootnote{metr.~emend.~from \emph{Viður at vígum} ‘id.’ \AM \emph{;} om.~\Regius}\Bfootnote{That this line is not an insertion by \AM\ is shown by the presence of \emph{Viðurr} in the corresponding part of the name-list in \Gylfaginning~20:11.}}.\eva

\bvb Grimner they called me at Garfrith’s \\
\ind but Gelding at Osmund’s, \\
but Keller when I pulled the sled, \\
\ind Throo at \inx[C]{Thing}[Things], \\
\ind Wither in battles.\evb\evg


\bvg\bva\mssnote{\Regius~10v/24, \AM~5v/10}%
\Ballnote{The names in this st.~appear in the list in \Gylfaginning~20:11 immediately following those of st.~49 in the order \emph{Ȯski, Ȯmi, Jafn·hǫ́r, Biflindi, Gǫndlir, Har·bárðr}.}%
\alst{Ȯ}ski ok \alst{Ȯ}mi, \hld\ \alst{Ja}fn·hǫ́r ok \edtext{Biflindi}{\Afootnote{\emph{Blindi} \RegiusProse}}, &
\ind \edtext{\alst{G}ǫndlir}{\Afootnote{\emph{Gęldnir} \Upsaliensis}} ok \edtrans{Hár·barðr með \alst{g}oðum}{Hoarbeard among Gods}{\Bfootnote{Viz.~in \textlink{Harbardsljod}, when conversing with the god Thunder.}}.\eva

\bvb Wisher and Omer, Evenhigh and Bivling, \\
\ind Gandler and Hoarbeard among Gods.\evb\evg


\bvg\bva\mssnote{\Regius~10v/25, \AM~5v/11}%
\Ballnote{The two Weden-names in this st.~appear in the list in \Gylfaginning~20:11 immediately following those of st.~51 in the order \emph{Sviðurr, Sviðrir} with no critical variants.}%
\alst{S}viðurr ok \alst{S}viðrir \hld\ es ek hét at \alst{S}økk·mïmis &
\ind ok dulða’k þann hinn \alst{a}ldna \alst{jǫ}tun &
þȧ’s \edtext{\alst{M}ið·vitnis}{\Afootnote{so \AM \emph{;} \emph{Mið·viðnis} \Regius}} vas’k \hld\ hins \alst{m}ę́ra burar &
\ind \alst{o}rðinn \alst{ęi}n-bani.\eva

\bvb Swither and Swithrer, as I was called at Sinkmimer’s, \\
\ind and I deceived that olden ettin, \\
when of Midwitner’s famous son \\
\ind I became the lone slayer.\evb\evg


\bvg\bva\mssnote{\Regius~10v/28, \AM~5v/13}%
\Ballnote{Weden now turns to King Garfrith.  He reproaches his former protegé and foresees his imminent death.}%
\alst{Ǫ}lr est Gęir·røðr, \hld\ hęfr þú \alst{o}f·drukkit; &
\alst{m}iklu est hnugginn, \hld\ es þú est \alst{m}ïnu \edtrans{gęngi}{support}{\Afootnote{so \Regius \emph{;} \emph{gęði} ‘mind’ \AM}}, &
\edtrans{\alst{ǫ}llum \alst{ęi}n-hęrjum}{all the Oneharriers}{\Bfootnote{It is important to note that Garfrith is not said to have lost the support (\emph{gęngi}) of the Oneharriers but rather the Oneharriers themselves, yet the sense is the same: by breaking the Odinic code of conduct befitting his station, Garfrith has lost Weden’s favour and thus been excluded from the community of oath-bound warriors, his adopted sons, the Oneharriers.  On the other hand, a righteous king could expect their truce, as was the case for Hathkin the Good according to \Hakonarmal, composed after his death.  Thus \Hakonarmal\ 16/1--2 sees the god \inx[P]{Bray} greeting Hathkin in Walhall with the words: \emph{ęin-hęrja grið \hld\ skalt allra hafa; / þigg þú at ǫ̇sum ǫl.} ‘the truce of all the Oneharriers shalt thou have.  Receive ale from the \inx[P]{Eese}!’}} \hld\ ok \edtrans{\alst{Ó}ðins hylli}{Weden’s holdness}{\Bfootnote{Weden’s grace or favour; cf.~st.~43/1 n.}}.\eva

\bvb Thou art worse for ale, Garfrith; thou hast drunk too much! \\
Much hast thou lost when thou hast [lost] My support, \\
all the \inx[P]{Oneharriers}, and Weden’s \inx[C]{holdness}.\evb\evg


\bvg\bva\mssnote{\Regius~10v/30, \AM~5v/15}%
\alst{F}jǫlð þér sagða’k, \hld\ ęn þú \alst{f}átt of mant, &
\ind of þik \alst{v}éla \edtext{\alst{v}inir}{\linenum{|2|||3}\lemma{vinir, mïns vinar ‘friends, my friend’}\Bfootnote{Weden stresses his friendship with Garfrith by repeating the word \emph{vinr} ‘friend’.  The followers of a god were his friends; see \textlink{Havamal}[157] n.}}; &
\edtext{\alst{m}ę́ki liggja \hld\ sé’k \alst{m}ïns vinar &
\ind allan ï \alst{d}ręyra \alst{d}rifinn.}{\lemma{mę́ki liggja \hld\ sé’k mïns vinar / allan ï dręyra drifinn. ‘A sword I see lying, in my friend’s / dreary blood all drenched.’}\Bfootnote{A prophetic statement.}}\eva

\bvb Many things have I told thee, but thou recallest few; \\
\ind ’tis friends that deal with thee! \\
A sword I see lying, in My friend’s \\
\ind dreary blood all drenched.\evb\evg


\bvg\bva\mssnote{\Regius~10v/31, \AM~5v/16}%
\edtrans{\alst{Ę}gg-móðan}{edge-wearied}{\Bfootnote{Euphemistic; “slain by a piercing blade.”  Also found in \textlink{Hamdismal}[30]/2a.}} val \hld\ nú mun \alst{Y}ggr hafa, &
\ind þitt vęit’k \alst{l}íf of \alst{l}iðit; &
\edtext{\alst{v}arar}{\Afootnote{\emph{†vvárro†} \AM}} ’ro \edtrans{dísir}{the Dises}{\Bfootnote{The Norns, fates, who have determined Garfrith’s hour of death.  Cf.~\textlink{Fafnismal}[11]--13, \textlink{Atlamal}[25]/4, \textlink{Hamdismal}[30].}}, \hld\ nú knátt \alst{Ó}ðin séa; &
\ind nálga⸗sk \alst{m}ik ef þú \alst{m}ęgir!\eva

\bvb An edge-wearied corpse will Ug now have; \\
\ind I know thy life to be past. \\
Wary are the \inx[P]{Dises}; now dost thou see Weden. \\
\ind Draw near Me, if thou mayst!\evb\evg


\bvg\bva\mssnote{\Regius~11r/2, \AM~5v/18}%
\Ballnote{The names in this st.~(other than \emph{Óðinn} ‘Weden’) appear in the name-list in \Gylfaginning~20:11 immediately following those of st.~50.  They have the archetypal order \emph{Yggr, Þundr, Vakr, Skilfingr, Vǫ́fuðr, Hropta-týr, Gautr, Jalkr}, although only the ms.~\Trajectinus\ keeps all names; \RegiusProse\ and \Wormianus\ omit.~\emph{Jalkr} while \Upsaliensis\ omits the first 6 names (\emph{Yggr--Hropta-týr} ‘Ug \dots\ Roft-Tew’) and therefore only has \emph{Gautr, Jalkr} ‘Geat, Gelding’.}%
\edtrans{\alst{Ó}ðinn nú hęiti’k}{Weden I am now called}{\Bfootnote{Having dropped all disguises, the guest is no longer Grimner but Weden himself; note that his earlier list of names did not contain this name.}}, \hld\ \alst{Y}ggr áðan hét’k, &
\ind hét⸗umk \alst{Þ}undr fyr \alst{þ}at, &
\alst{V}akr ok Skilfingr, \hld\ \edtext{\alst{V}ǫ́fuðr}{\Afootnote{so \Regius\AM\RegiusProse \emph{;} \emph{Ná·fǫðr} (\emph{‘nafodr’}) \Trajectinus \emph{;} \emph{Vá·fǫðr} \Wormianus \emph{;} om.~\Upsaliensis}} ok Hropta-týr &
\ind \alst{G}autr ok Jalkr með \alst{g}oðum.\eva

\bvb Weden am I now called; Ug was I called earlier; \\
\ind I called myself Thound before that; \\
Wacker and Shilving, Waved and Roft-Tew, \\
\ind Geat and Gelding among the Gods.\evb\evg


\bvg\bva\mssnote{\Regius~11r/4, \AM~5v/20}%
\Ballnote{No trace of this stanza is found in the list of names in \Gylfaginning~20:11 which suggests that it may be a later addition common to \Regius\ and \AM.  It is also possible that these names were left out since Snorre considered them to be serpent-names on the basis of st.~35 above, which is quoted in \Gylfaginning~16:9.}%
\edtrans{\alst{Ó}fnir ok Sváfnir}{Ovner and Swebner}{\Bfootnote{Both names are listed among poetic names of Weden in Þul \emph{Óðins} (in \Skp\ 3).  They are also the names of two serpents in 35/3 above.}} \hld\ \edtext{es}{\Afootnote{om.~\AM}} hygg’k at \alst{o}rðnir sé &
\ind \alst{a}llir \edtext{af}{\Afootnote{\emph{at} \AM}} \alst{ęi}num mér.“\eva

\bvb Ovner and Swebner, whom, I ween, have come \\
\ind all from Me alone.”\evb\evg


\bpg\bpa[3]\mssnote{\Regius~11r/5, \AM~5v/21}%
Geir·røðr konungr sat, ok hafði sverð um kné sér ok brugðit til miðs. En er hann heyrði, at Óðinn var þar \edtext{kominn,}{\Afootnote{add.~\emph{þá} \AM}} stóð hann upp, ok vildi taka Óðin frá eldi’num. Sverð’it slapp ór hendi \edtext{hǫ́num,}{\Afootnote{add.~\emph{ok} \AM}}  vissu hjǫlt’in niðr. Konungr drap fǿti, ok steypti⸗sk á-fram, en sverð’it stóð í gǫgnum hann, ok fekk \edtext{hann}{\Afootnote{\emph{þar af} \AM}} bana. \edtext{Óðinn hvarf þá.}{\Afootnote{om. \AM}}
En Agnarr \edtext{var þar}{\Afootnote{\emph{varð} \AM}} konungr \edtext{lengi síðan.}{\Afootnote{om. \AM}}\epa

\bpb King Garfrith sat and had a sword about his knee, and it was brandished half-way up.  And when he heard that Weden had come there, he stood up and would take Weden from the fire.  The sword slipped out of his hand; the hilt pointed downwards.  The king lost his footing and fell forward, but the sword went through him and he received his bane.  Weden then disappeared, but Ayner was king there for a long while afterwards.\epb\epg

\section{Appendix}

\subsection{Origin of the Nation of the Longbeards (\emph{Origo Gentis Langbardorum})}

The following excerpt is the first chapter of the 7th century text dealing with the legendary prehistory of the Longbeards.  It is edited here due to its close resemblance to P2 of the prologue above and since certain name-forms require discussion in a way that would not be entirely in place in a footnote.

\bpg\bpa[0]%
\pva Est insula qui dicitur \edtrans{Scadanan, quod interpretatur excidia}{Shedeney, which is interpreted as “destruction”}{\Bfootnote{\emph{Scadanan} ‘Shedeney’ is the same word as Scandinavia or Scania, \emph{Scanza} in Jordanes.  The word derives from PGmc *\emph{Skadin-awjō} ‘Shede’s island’, where Shede is probably the goddess (ON \emph{Skaði}) known from Norse sources. --- The gloss reveals that the author was not unfamiliar with Germanic idiom; cf.~OHG \emph{skado} ‘damage’.}}, in partibus aquilonis, ubi multae gentes habitant; inter quos erat gens parva quae Winnilis vocabatur.
\pva Et erat cum eis mulier nomine Gambara, habebatque duos filios, nomen uni Ybor et nomen alteri agio; ipsi cum matre sua nomine Gambara principatum tenebant super Winniles.
\pva Moverunt se ergo duces Wandalorum, id est Ambri et Assi, cum exercitu suo,
\pva et dicebant ad Winniles: “Aut solvite nobis tributa, aut praeparate vos ad pugnam et pugnate nobiscum”.
\pva Tunc responderunt Ybor et Agio cum matre sua Gambara: “Melius est nobis pugnam praeparare, quam wandalis tributa persolvere”.
\pva Tunc Ambri et Assi, hoc est duces Wandalorum, rogaverunt \edtrans{Godan}{Weden}{\Bfootnote{A form derived from earlier \emph{Wódan} (attested on the runic Nordendorf fibula) through the sound change \emph{w-} > \emph{g-}.}}, ut daret eis super Winniles victoriam.
\pva Respondit Godan dicens: “Quos sol surgente antea videro, ipsis dabo victoriam”.
\pva Eo tempore Gambara cum duobus filiis suis, id est Ybor et Agio, rogaverunt \edtrans{Fream}{Frie}{\Bfootnote{Weden’s wife, Frie.  Acc.~sg.~of \emph{Frea}, cognate with OHG \emph{Frija} (\textlink{MerseburgTwo} 4), ON \emph{Frigg}, both from PGmc.~*\emph{Frijō}.
Frie is not to be confused with the different goddess Frow (ON \emph{Fręyja}, whose name, though superficially similar to \emph{Frea}, in fact comes from PGmc.~*\emph{frawjǭ} ‘lady’, corresponding to OHG \emph{frouwa}.}}, uxorem Godam, ut ad Winniles esset propitia.
\pva Tunc Frea dedit consilium, ut sol surgente venirent Winniles et mulieres eorum crines solutae circa faciem in similitudinem barbae et cum viris suis venirent.
\pva Tunc luciscente sol dum surgeret, giravit Frea, uxor Godan, lectum ubi recumbebat vir eius, et fecit faciem eius contra orientem, et excitavit eum.
\pva Et ille aspiciens vidit Winniles et mulieres ipsorum habentes crines solutas circa faciem;
et ait: “Qui sunt isti longibarbae?”
\pva Et dixit Frea ad Godan: “\edtrans{Sicut dedisti nomen, da illis et victoriam}{As thou hast given the name, give them also the victory}{\Bfootnote{Likely an early attestation of the custom of the name-gift (ON \emph{gefa at nafn-fęstr} ‘give for the fastening of a name’) found in numerous later Norse Saws.
Cf.~\Skaldskaparmal\ 53: \emph{Þá svarar konungr: „Þú, sveinn, hefir gefit mér nafn, at ek skal heita Hrólfr kraki, en þat er títt, at gjǫf skal fylgja nafn-festi. [...]“} ‘Then answers the king: “Thou, swain, hast given me a name, that I shall be called Rotholf the pole; yet it is customary that a gift should follow the fastening of a name. [...]”’}}”.
\pva Et dedit eis victoriam, ut ubi visum esset vindicarent se et victoriam haberent.
\pva Ab illo tempore Winnilis Langobardi vocati sunt.\epa

\bpb\pvb There is an island which is called Shedeney, which is interpreted as ‘destruction’, in the parts of the north, where many nations dwell; among whom was a small nation which was called \emph{Winnilis}.
\pvb And there was with them a woman by name \emph{Gambara}, and she had two sons, the name of one was \emph{Ybor} and the name of the other \emph{Agio}; they with their mother by name Gambara held leadership over the Winnilis.
\pvb Then the princes of the \emph{Wandali}, that is \emph{Ambri} and \emph{Assi}, moved with their host,
\pvb and said to the Winnilis: “Either pay tributes to us, or prepare yourselves for battle and fight with us.”
\pvb Then Ybor and Agio with their mother Gambara answered: “It is better for us to prepare for battle than to pay tributes to the Wandali.”
\pvb Then Ambri and Assi, that is the princes of the Wandali, asked Weden that he might give to them victory over the Winnilis.
\pvb Weden answered, saying: “Those whom I shall see first at the rising sun, to them I will give the victory.”
\pvb At that time Gambara with her two sons, that is Ybor and Agio, asked Frie, the wife of Weden, that she might be favorable toward the Winnilis.
\pvb Then Frie gave counsel, that at the rising sun the Winnilis should come and their women with hair let down around the face in the likeness of a beard, and they should come with their men.
\pvb Then as it grew light while the sun was rising, Frie, wife of Weden, turned the bed where her husband was reclining, and made his face turn toward the east, and awakened him.
\pvb And he, looking out, saw the Winnilis and their women having hair let down around the face;
and he said: “Who are these long-beards?”
\pvb And Frie said to Weden: “As thou hast given the name, give them also the victory.”
\pvb And he gave them victory, so that wherever it seemed [fit] they might avenge themselves and have victory.
\pvb From that time the Winnilis are called \emph{Langobardi}.\epb\epg
